% !TeX root = ../../Book4.tex
% 纲要标题：## 第3章　可表函子、Yoneda 引理与伴随
% 纲要位置：计划纲要.md，第 3157 行。

\chapter{可表函子、Yoneda 引理与伴随}
\label{b4:ch:03}
\BookChapterNumber{b4:ch:03}

\begin{chapterabstract}
本章以可表函子和 Yoneda 引理刻画对象与自然变换，建立伴随的 Hom 双射、单位与余单位两种表述，并证明左伴随保持余极限、右伴随保持极限。自由模与张量积提供具体例子，最后从伴随构造单子和余单子。
\end{chapterabstract}

\begin{learninggoals}
\begin{itemize}
\item 为一个具体集合值函子找出表示对象、泛元素和自然双射，并辨认协变与反变的方向。
\item 证明 Yoneda 引理，分别检查对象变量与函子变量中的自然性，利用它求自然变换。
\item 从伴随双射构造单位与余单位，也能从三角恒等式恢复双射及其两变量自然性。
\item 写出自由—忘却、交换化以及 Hom–张量伴随的具体映射，核对非交换模的侧别。
\item 用泛性质证明伴随保持相应的极限，给出不能保持另一类极限的反例，并用形式运算解释单子公理。
\end{itemize}
\end{learninggoals}

设 \(R\) 为含幺环，\(M\) 为幺性左 \(R\)-模。给出 \(R\rightarrow M\) 的模同态，只需说出 \(1\in R\) 的像：若此像为 \(m\)，同态必将 \(r\) 送到 \(rm\)。于是 \(M\) 的元素可以通过从同一个对象 \(R\) 出发的态射来辨认。若 \(N\) 是另一个幺性左 \(R\)-模，而 \(f:M\rightarrow N\) 是模同态，映射 \(m\mapsto f(m)\) 又恰好对应同态 \(R\rightarrow M\xrightarrow{f}N\) 的复合。这里值得追问的，不只是某两个集合之间碰巧有双射，而是这些双射能否随着目标对象变化而保持相容。本章从这个问题进入可表性、Yoneda 引理与伴随。

\ProseParagraphBreak
本章使用\chapref{b4:ch:01}的函子、自然变换及全忠实性，以及\chapref{b4:ch:02}的极限、余极限。模论例子调用第二卷第3.1节的自由模泛性质、第6.1节的张量积构造和第6.3节的 Hom–张量伴随；这些构造已在第二卷证明，本章研究它们所共有的范畴性质。交换化的例子另用群的交换子和商群。

\ProseParagraphBreak
可对照 Weibel《An Introduction to Homological Algebra》中的 Yoneda 嵌入\cite[Example A.3.4]{Weibel1994HomologicalAlgebra}及伴随的表述\cite[A.6.1; Theorem A.6.2]{Weibel1994HomologicalAlgebra}。本章依次证明泛元素判据、Yoneda 引理、伴随的两种刻画及保持定理；自由模的生成元和张量积的纯张量为这些一般公式提供了具体检验。


% !TeX root = ../../Book4.tex
% 纲要标题：### 3.1 可表函子
% 纲要位置：计划纲要.md，第 3159 行。

\section{可表函子}
\label{b4:sec:0301}

自由对象把指定元素的问题化为映射问题，张量积则把平衡双可加映射化为群同态。可表性把这些对应写成对目标或源对象自然的 Hom 描述。

% 正文按纲要写入；各节的基础结果在本章证明。

% 纲要内容（仅作写作计划，尚非教材正文）：
%
% * Hom 函子及表示对象
% * 张量积和自由对象的可表性解释；非交换环上的平衡双可加映射与群同态
% * 表示对象的典范唯一性
%

\subsection{Hom 函子及表示对象}
\label{b4:subsec:0301:01}

自由模的泛性质把一组任意指定的元素变成模同态；张量积的泛性质把平衡双可加映射变成群同态。这两种对应都随着目标对象的改变而相容。本节把这种情形统一为一个问题：一个取值于集合的函子，能否用从某个固定对象出发或到达某个固定对象的态射来描述？

\ProseParagraphBreak
设 \(\mathcal C\) 为局部小范畴，\(A\in\mathcal C\)。记
\[
h^A(X)=\operatorname{Hom}_{\mathcal C}(A,X),\qquad
h_A(X)=\operatorname{Hom}_{\mathcal C}(X,A).
\]
对态射 \(u:X\rightarrow Y\)，分别令 \(h^A(u)(f)=u\circ f\) 与 \(h_A(u)(g)=g\circ u\)。复合的结合律保证 \(h^A:\mathcal C\rightarrow\mathbf{Set}\) 是协变函子，\(h_A:\mathcal C^{\mathrm{op}}\rightarrow\mathbf{Set}\) 是反变函子。下标固定的是目标，上标固定的是源；只有自由变动的那个变量决定函子的变性。
\symidx{\(h^A,h_A\)：以 \(A\) 为固定源或固定目标的 Hom 函子}

\begin{definition}\label{b4:def:representable-functor}
设 \(\mathcal C\) 为局部小范畴，对 \(A\in\mathcal C\) 记 \(h^A=\operatorname{Hom}_{\mathcal C}(A,-)\)、\(h_A=\operatorname{Hom}_{\mathcal C}(-,A)\)。给定函子 \(F:\mathcal C\rightarrow\mathbf{Set}\)。若存在对象 \(A\in\mathcal C\) 和自然同构 \(\theta:h^A\xrightarrow{\sim}F\)，则称 \(F\) 为\term[kebiao-hanzi]{可表函子}[representable functor]，称 \((A,\theta)\) 为 \(F\) 的一个表示，\(A\) 为表示对象。对于 \(F:\mathcal C^{\mathrm{op}}\rightarrow\mathbf{Set}\)，相应要求为自然同构 \(\theta:h_A\xrightarrow{\sim}F\)。
\end{definition}

这里既要求每个 \(\theta_X\) 为双射，也要求这些双射组成自然变换。协变情形的相容条件为
\[
F(u)\bigl(\theta_X(f)\bigr)=\theta_Y(u\circ f)
\quad(f:A\rightarrow X,\ u:X\rightarrow Y).
\]
只知道各个集合具有相同基数，并未给出一个表示。把同一个对象在不同目标中的作用拼在一起，正是自然性所承担的工作。

\begin{example}\label{b4:exm:representable-forgetful}
设 \(R\) 为含幺环，\(U:R\text{-}\mathbf{Mod}\rightarrow\mathbf{Set}\) 为左模的忘却函子。左正则模 \({}_RR\) 表示 \(U\)，相应双射为
\[
\operatorname{Hom}_R(R,M)\longrightarrow U(M),\qquad f\longmapsto f(1).
\]
逆映射把 \(m\in M\) 送到 \(r\mapsto rm\)，故无需选择 \(M\) 的基。若 \(v:M\rightarrow N\) 左 \(R\) 线性，则 \((v\circ f)(1)=v(f(1))\)，这恰好是自然性。特别地，\(\mathbb Z\) 表示交换群的底集合函子。
\end{example}

可表性也包括反变例子。对集合 \(X\)，把子集 \(B\subseteq X\) 对应到特征函数 \(\chi_B:X\rightarrow\{0,1\}\)。等式 \(\chi_{u^{-1}(B)}=\chi_B\circ u\) 表明，取逆像的幂集函子由二元集合表示。这里必须取逆像；取直接像得到另一个协变函子，不能与这一表示混用。

\subsection{张量积和自由对象的可表性解释}
\label{b4:subsec:0301:02}

固定一个集合 \(X\) 和一个含幺环 \(R\)。函子
\[
M\longmapsto U(M)^X
\]
把左 \(R\) 模 \(M\) 送到从 \(X\) 到其底集合的全部映射；模同态通过后复合作用。第二卷第3.1节的自由模泛性质已说明，自由左模 \(R^{(X)}=\bigoplus_{x\in X}Re_x\) 满足
\[
\operatorname{Hom}_R(R^{(X)},M)\xrightarrow{\sim}U(M)^X,
\qquad f\longmapsto\bigl(x\mapsto f(e_x)\bigr).
\]
逆映射将有限和 \(\sum_x r_xe_x\) 送到 \(\sum_xr_x\cdot a(x)\)。有限支集保证这个公式在一般模中有意义；对模同态 \(v\)，等式 \(v(\sum_xr_x\cdot a(x))=\sum_xr_x\cdot v(a(x))\) 给出自然性。因此自由模表示的是“指定 \(X\) 个元素”的函子。

\ProseParagraphBreak
再固定右 \(R\) 模 \(M_R\) 与左 \(R\) 模 \({}_RN\)。对交换群 \(A\)，记 \(\operatorname{Bal}_R(M,N;A)\) 为所有双可加且满足
\[
b(mr,n)=b(m,rn)
\]
的映射 \(b:M\times N\rightarrow A\) 组成的集合。交换群同态 \(v:A\rightarrow A'\) 将 \(b\) 送到 \(v\circ b\)，于是得到 \(\mathbf{Ab}\rightarrow\mathbf{Set}\) 的函子。第二卷第6.1节的张量积构造及泛性质给出表示
\[
\operatorname{Hom}_{\mathbb Z}(M\otimes_RN,A)
\xrightarrow{\sim}\operatorname{Bal}_R(M,N;A),
\qquad f\longmapsto\bigl((m,n)\mapsto f(m\otimes n)\bigr).
\]
关于 \(A\) 的自然性就是 \(v(f(m\otimes n))=(v\circ f)(m\otimes n)\)。这里不假定 \(R\) 交换，张量积首先是交换群；若还要把它看成某个环上的模，必须另给相容的双模结构。当 \(R\) 交换且 \(M,N,A\) 都是 \(R\) 模时，\(M\otimes_RN\) 具有自然的 \(R\) 模结构，\(R\) 双线性映射便对应 \(R\) 模同态；这才是通常所说的“双线性映射化为线性映射”。

\[
\begin{tikzcd}[column sep=large,row sep=large]
M\times N\arrow[r,"{(m,n)\mapsto m\otimes n}"]\arrow[dr,"b"']&
M\otimes_RN\arrow[d,dashed,"\substack{\exists!\,f\\\text{群同态}}" description]\\
&A.
\end{tikzcd}
\]

这两例不仅给出表示对象，也各自带有一个特别的元素：自由模带有映射 \(x\mapsto e_x\)，张量积带有平衡映射 \((m,n)\mapsto m\otimes n\)。每个所研究的映射都由这个元素经唯一态射得到。下面的命题把这一观察写成可直接使用的判据。

\begin{proposition}[泛元素判据]\label{b4:prop:universal-element}
设 \(\mathcal C\) 为局部小范畴，\(F:\mathcal C\rightarrow\mathbf{Set}\) 为函子。给定 \(A\in\mathcal C\) 及 \(u\in F(A)\)，记 \(h^A=\operatorname{Hom}_{\mathcal C}(A,-)\)。公式
\[
\theta_X(f)=F(f)(u),\qquad f:A\rightarrow X,
\]
定义自然变换 \(h^A\Rightarrow F\)。它是自然同构，当且仅当每个 \(X\in\mathcal C\) 及 \(x\in F(X)\) 都对应唯一态射 \(f:A\rightarrow X\)，使 \(F(f)(u)=x\)。此时称 \(u\) 为该表示的\term[fanyuan-su]{泛元素}[universal element]。每个表示 \(\theta:h^A\xrightarrow{\sim}F\) 都由 \(u=\theta_A(1_A)\) 按上述公式得到。
\end{proposition}
\begin{proof}
若 \(g:X\rightarrow Y\)，则
\[
F(g)\bigl(\theta_X(f)\bigr)=F(g)F(f)(u)
=F(gf)(u)=\theta_Y(gf),
\]
故 \(\theta\) 自然。每个分量为双射，等价于所述存在唯一性。反过来，若 \(\theta\) 已是一个表示，则将自然性用于 \(f:A\rightarrow X\) 以及 \(1_A\)，得到
\[
\theta_X(f)=\theta_X(f\circ1_A)=F(f)\bigl(\theta_A(1_A)\bigr).
\]
所以这个泛元素给出的公式恢复全部分量。
\end{proof}

对反变函子 \(F:\mathcal C^{\mathrm{op}}\rightarrow\mathbf{Set}\)，相同论证给出 \(u\in F(A)\) 与双射 \(f:X\rightarrow A\mapsto F(f)(u)\)。自由模例子的泛元素是一个以 \(X\) 为指标的元素族，张量积例子的泛元素是一个平衡映射；“元素”指 \(F(A)\) 中的成员，并不一定是对象 \(A\) 底集合中的一个点。

\subsection{表示对象的典范唯一性}
\label{b4:subsec:0301:03}

\begin{proposition}\label{b4:prop:representing-unique}
设 \(\mathcal C\) 为局部小范畴，\(F:\mathcal C\rightarrow\mathbf{Set}\) 为函子，\(A,B\in\mathcal C\)，\(u\in F(A)\)、\(v\in F(B)\) 均为 \(F\) 的泛元素。则有唯一同构 \(s:A\rightarrow B\)，使 \(F(s)(u)=v\)。对于反变函子 \(F:\mathcal C^{\mathrm{op}}\rightarrow\mathbf{Set}\) 及其泛元素 \(u\in F(A)\)、\(v\in F(B)\)，亦有唯一同构 \(s:A\rightarrow B\)，使 \(F(s)(v)=u\)。
\end{proposition}
\begin{proof}
先处理协变情形。由 \((A,u)\) 的泛性质，有唯一 \(s:A\rightarrow B\) 满足 \(F(s)(u)=v\)；由 \((B,v)\) 的泛性质，有唯一 \(t:B\rightarrow A\) 满足 \(F(t)(v)=u\)。于是
\[
F(ts)(u)=u=F(1_A)(u),\qquad
F(st)(v)=v=F(1_B)(v).
\]
泛性质的唯一性分别给出 \(ts=1_A\) 与 \(st=1_B\)，故 \(s\) 是同构。任一满足条件的同构首先也是满足条件的态射，已由构造中的唯一性确定。反变情形对 \(\mathcal C^{\mathrm{op}}\) 应用协变结论，再把所得同构取逆，便得到所写的方向和等式。
\end{proof}

“唯一”指保持表示数据的同构唯一，并非两个底对象之间只有一个同构。例如 \(\mathbb Z\) 作为交换群有恒等与取负两个自同构，但保持代表底集合函子所用泛元素 \(1\) 的只有恒等。

\ProseParagraphBreak
对含幺环 \(R\) 上的右模 \(M_R\) 与左模 \({}_RN\)，设交换群 \(T,T'\) 分别连同平衡双可加映射 \(\tau:M\times N\rightarrow T\)、\(\tau':M\times N\rightarrow T'\)，都满足张量积的泛性质。这两个映射就是各自表示的泛元素，因此存在唯一群同构 \(s:T\rightarrow T'\)，使
\[
s\bigl(\tau(m,n)\bigr)=\tau'(m,n)
\qquad(m\in M,\ n\in N).
\]
任意一个群同构未必保持这两组表示数据。

\ProseParagraphBreak
自然性还会决定表示对象之间的态射。前一命题只讨论同构，下一节将证明更强的结论：Hom 函子之间的任意自然变换，都来自唯一的对象间态射。表示对象的唯一性由此成为一个一般事实的特例。

% !TeX root = ../../Book4.tex
% 纲要标题：### 3.2 Yoneda 引理
% 纲要位置：计划纲要.md，第 3165 行。

\section{Yoneda 引理}
\label{b4:sec:0302}

一个对象接受的全部映射，是否已经决定这个对象及其间的态射？Yoneda 引理通过在恒等态射处求值回答这个问题，并给出所需自然性。

% 正文按纲要写入；各节的基础结果在本章证明。

% 纲要内容（仅作写作计划，尚非教材正文）：
%
% * Yoneda 嵌入；小范畴中的全忠实函子与局部小范畴的 Hom 双射
% * 引理的证明及自然性
% * 用 Hom 函子探测对象与态射
%

\subsection{Yoneda 嵌入}
\label{b4:subsec:0302:01}

对局部小范畴 \(\mathcal C\) 中的对象 \(A\)，函子 \(h_A=\operatorname{Hom}_{\mathcal C}(-,A)\) 记录从每个对象到 \(A\) 的态射。若 \(a:A\rightarrow B\)，后复合给出自然变换
\[
h_a:h_A\Longrightarrow h_B,
\qquad (h_a)_X(f)=a\circ f.
\]
自然性来自 \(a\circ(f\circ u)=(a\circ f)\circ u\)，其中 \(u:X'\rightarrow X\)。又有 \(h_{1_A}=1_{h_A}\) 及 \(h_{ba}=h_b\circ h_a\)。当 \(\mathcal C\) 小时，因此得到函子
\[
\mathsf y:\mathcal C\longrightarrow\operatorname{Fun}(\mathcal C^{\mathrm{op}},\mathbf{Set}),
\qquad A\longmapsto h_A,\quad a\longmapsto h_a.
\]
它称为\term[yoneda-qianru]{Yoneda 嵌入}[Yoneda embedding]。名称中的“嵌入”将在下面证明为全忠实性；一般不能把它理解成对象集合上的字面包含。
\symidx{\(\mathsf y\)：Yoneda 嵌入}

\ProseParagraphBreak
这里选用 \(h_A\)，是因为 \(a:A\rightarrow B\) 使 \(h_A\rightarrow h_B\)，在表示对象变量上同向变化。若改用 \(h^A=\operatorname{Hom}_{\mathcal C}(A,-)\)，预合成 \(a\) 得到的却是 \(h^B\rightarrow h^A\)，因此在小范畴情形得到的是 \(\mathcal C^{\mathrm{op}}\rightarrow\operatorname{Fun}(\mathcal C,\mathbf{Set})\)。

\ProseParagraphBreak
若 \(\mathcal C\) 仅局部小，下文按分量处理 \(h_A\) 以及涉及它的自然变换，不预设整个集合值函子范畴都在\chapref{b4:ch:01}的大小范围内。引理将证明，以 \(h_A\) 为源、\(F\) 为目标的自然变换由集合 \(F(A)\) 参数化，因而可表函子之间的自然变换构成集合。这些大小约定保证后面所写的 \(\operatorname{Nat}(h_A,F)\) 确实是集合；双射本身的构造则由恒等态射与自然性给出。对于仅局部小的范畴，本书所说的“Yoneda 全忠实性”指下文给出的 Hom 集与自然变换集合之间的双射。

\ProseParagraphBreak
在集合范畴，\(h_A(X)\) 是所有 \(X\) 指标的 \(A\) 中元素族；改变指标的映射 \(X'\rightarrow X\)，就是把元素族拉回。一个映射 \(a:A\rightarrow B\) 对每个族逐项应用 \(a\)，于是自动产生相容的族变换。对于一般范畴，态射 \(X\rightarrow A\) 取代了元素族。关键问题是：每一种对所有 \(X\) 都相容的变换，是否都这样得到？

\subsection{Yoneda 引理的证明与自然性}
\label{b4:subsec:0302:02}

对源、目标相同的两个函子 \(H,F\)，记 \(\operatorname{Nat}(H,F)\) 为从 \(H\) 到 \(F\) 的自然变换类。\term[yoneda-yinli]{Yoneda 引理}[Yoneda lemma]说明，以可表函子为源时，这个类是集合，而且只需计算一个分量中的一个值。

\begin{theorem}[Yoneda 引理]\label{b4:thm:yoneda}
设 \(\mathcal C\) 为局部小范畴，\(F:\mathcal C^{\mathrm{op}}\rightarrow\mathbf{Set}\) 为函子，\(A\in\mathcal C\)，记 \(h_A=\operatorname{Hom}_{\mathcal C}(-,A)\)。求值映射
\[
\Phi_{A,F}:\operatorname{Nat}(h_A,F)\longrightarrow F(A),
\qquad \alpha\longmapsto\alpha_A(1_A)
\]
是双射，因而自然变换类 \(\operatorname{Nat}(h_A,F)\) 是集合。其逆把 \(u\in F(A)\) 送到自然变换 \(\alpha^u\)，其中
\[
\alpha^u_X(f)=F(f)(u)\qquad(f:X\rightarrow A).
\]
这一双射在 \(A\) 与 \(F\) 上自然：设 \(a:A\rightarrow B\) 为 \(\mathcal C\) 中的态射，\(h_B=\operatorname{Hom}_{\mathcal C}(-,B)\)，\(h_a:h_A\Rightarrow h_B\) 为后复合 \(a\) 得到的自然变换。改变表示对象时，自然变换一边预合成 \(h_a\)，函子值一边通过 \(F(a)\) 变化。改变函子时，自然变换一边后复合，函子值一边作用相应的 \(A\) 分量。具体地，若 \(G:\mathcal C^{\mathrm{op}}\rightarrow\mathbf{Set}\) 为函子，\(\alpha:h_B\Rightarrow F\)、\(\beta:F\Rightarrow G\) 为自然变换，则
\[
\begin{aligned}
\Phi_{A,F}(\alpha\circ h_a)&=F(a)\bigl(\Phi_{B,F}(\alpha)\bigr),\\
\Phi_{A,G}(\beta\circ\gamma)&=\beta_A\bigl(\Phi_{A,F}(\gamma)\bigr)
\qquad(\gamma:h_A\Rightarrow F).
\end{aligned}
\]
对协变函子 \(F:\mathcal C\rightarrow\mathbf{Set}\)，记 \(h^A=\operatorname{Hom}_{\mathcal C}(A,-)\)，相应结论为
\(\operatorname{Nat}(h^A,F)\cong F(A)\)，逆映射仍由 \(f:A\rightarrow X\mapsto F(f)(u)\) 给出。
\end{theorem}
\begin{proof}
给定 \(u\in F(A)\)，首先验证公式确实定义自然变换。对 \(v:X'\rightarrow X\) 及 \(f:X\rightarrow A\)，反变函子公理给出
\[
\alpha^u_{X'}(fv)=F(fv)(u)=F(v)F(f)(u)
=F(v)\bigl(\alpha^u_X(f)\bigr).
\]
所以 \(\alpha^u\) 自然，且 \(\alpha^u_A(1_A)=u\)。反过来，若 \(\alpha:h_A\Rightarrow F\) 自然，任意 \(f:X\rightarrow A\) 都满足 \(h_A(f)(1_A)=1_A\circ f=f\)。所以所有这样的 \(f\) 都由同一个 \(1_A\) 经前复合得到。将自然性用于这条箭头和这个元素，得到
\[
\alpha_X(f)=F(f)\bigl(\alpha_A(1_A)\bigr).
\]
\[
\begin{tikzcd}[column sep=large,row sep=large]
h_A(A)\arrow[r,"\alpha_A"]\arrow[d,"h_A(f)"']&F(A)\arrow[d,"F(f)"]\\
h_A(X)\arrow[r,"\alpha_X"']&F(X).
\end{tikzcd}
\]
这说明从 \(\alpha_A(1_A)\) 按所给公式重建的自然变换就是 \(\alpha\)，故两个映射互逆。所有自然变换恰为由 \(u\in F(A)\) 构造出的 \(\alpha^u\)，所以它们构成一个由集合 \(F(A)\) 参数化的集合。

现在检查两个参数的自然性。对 \(a:A\rightarrow B\) 及 \(\alpha:h_B\Rightarrow F\)，有
\[
(\alpha\circ h_a)_A(1_A)=\alpha_A(a)
=F(a)\bigl(\alpha_B(1_B)\bigr).
\]
由 \(\Phi\) 在恒等态射处求值的定义，这正是下列方块在任意 \(\alpha\) 上的交换条件：
\[
\begin{tikzcd}[column sep=large,row sep=large]
\operatorname{Nat}(h_B,F)\arrow[r,"\Phi_{B,F}"]\arrow[d,"{(-)\circ h_a}"']&F(B)\arrow[d,"F(a)"]\\
\operatorname{Nat}(h_A,F)\arrow[r,"\Phi_{A,F}"']&F(A).
\end{tikzcd}
\]
对 \(\beta:F\Rightarrow G\) 及 \(\gamma:h_A\Rightarrow F\)，则
\[
(\beta\circ\gamma)_A(1_A)
=\beta_A\bigl(\gamma_A(1_A)\bigr).
\]
这同样说明后复合 \(\beta\) 与在函子值上作用 \(\beta_A\) 相容，即
\[
\begin{tikzcd}[column sep=large,row sep=large]
\operatorname{Nat}(h_A,F)\arrow[r,"\Phi_{A,F}"]\arrow[d,"{\beta\circ(-)}"']&F(A)\arrow[d,"\beta_A"]\\
\operatorname{Nat}(h_A,G)\arrow[r,"\Phi_{A,G}"']&G(A).
\end{tikzcd}
\]
协变形式对相反范畴 \(\mathcal C^{\mathrm{op}}\) 应用刚证的结论即可；在原范畴中，所用态射方向恰变成 \(A\rightarrow X\)。
\end{proof}

自然性迫使一个变换在所有态射上的值由它在表示对象的恒等态射上的值决定，函子公理则保证得到的公式在所有对象间相容。这就是前一节泛元素判据的公式为何适用于每个自然变换，而不只适用于自然同构。

\begin{corollary}\label{b4:cor:yoneda-embedding}
设 \(\mathcal C\) 为局部小范畴。对任意 \(A,B\in\mathcal C\)，记 \(h_A=\operatorname{Hom}_{\mathcal C}(-,A)\)、\(h_B=\operatorname{Hom}_{\mathcal C}(-,B)\)，以 \(h_a:h_A\Rightarrow h_B\) 表示后复合态射 \(a:A\rightarrow B\) 得到的自然变换。映射
\[
\operatorname{Hom}_{\mathcal C}(A,B)\longrightarrow\operatorname{Nat}(h_A,h_B),
\qquad a\longmapsto h_a
\]
是双射。若 \(\mathcal C\) 进一步为小范畴，则 Yoneda 嵌入
\[
\mathsf y:\mathcal C\longrightarrow\operatorname{Fun}(\mathcal C^{\mathrm{op}},\mathbf{Set}),
\qquad A\longmapsto h_A,\quad a\longmapsto h_a
\]
为全忠实函子。对于局部小范畴 \(\mathcal C\)，若 \(h_A\) 与 \(h_B\) 自然同构，则 \(A\cong B\)。
\end{corollary}
\begin{proof}
在 Yoneda 引理中取 \(F=h_B\)，右边为 \(h_B(A)=\operatorname{Hom}_{\mathcal C}(A,B)\)，逆对应正是后复合。因此态射映射为双射。当 \(\mathcal C\) 小时，函子 \(\mathsf y\) 由这些映射给出，因而全且忠实。若自然同构 \(\alpha:h_A\xrightarrow{\sim}h_B\) 及其逆分别来自 \(a:A\rightarrow B\)、\(b:B\rightarrow A\)，则
\[
h_{ba}=\alpha^{-1}\alpha=1_{h_A}=h_{1_A},\qquad
h_{ab}=1_{h_B}=h_{1_B}.
\]
上面双射的单射性给出 \(ba=1_A\)、\(ab=1_B\)，故 \(a\) 是同构。
\end{proof}

这一嵌入的标准表述可参阅 Weibel《An Introduction to Homological Algebra》\cite[Example A.3.4]{Weibel1994HomologicalAlgebra}。对于协变 Hom 函子，态射 \(s:B\rightarrow A\) 对应自然变换 \(h^A\Rightarrow h^B\)，其分量为
\[
\operatorname{Hom}_{\mathcal C}(A,X)\longrightarrow
\operatorname{Hom}_{\mathcal C}(B,X),\qquad f\longmapsto f\circ s.
\]
因此 \(\operatorname{Nat}(h^A,h^B)\cong\operatorname{Hom}_{\mathcal C}(B,A)\)。固定源的 Hom 函子在表示对象上反向变化，正是由于这里要预合成 \(s\)。

\subsection{用 Hom 函子探测对象与态射}
\label{b4:subsec:0302:03}

\begin{proposition}\label{b4:prop:hom-detects-morphisms}
设 \(\mathcal C\) 为局部小范畴，\(f:A\rightarrow B\) 为态射。则：\(f\) 为单态射，当且仅当每个 \(X\in\mathcal C\) 的后复合映射 \(\operatorname{Hom}_{\mathcal C}(X,A)\rightarrow\operatorname{Hom}_{\mathcal C}(X,B)\) 都单射；\(f\) 为满态射，当且仅当每个 \(X\) 的前复合映射 \(\operatorname{Hom}_{\mathcal C}(B,X)\rightarrow\operatorname{Hom}_{\mathcal C}(A,X)\) 都单射。此外，\(f\) 为同构，当且仅当上述后复合映射对每个 \(X\) 都双射。
\end{proposition}
\begin{proof}
前两项分别把单态射与满态射的消去条件按 Hom 集重写，没有增加条件。若 \(f\) 是同构，后复合 \(f^{-1}\) 给出每个映射的逆。反过来，取 \(X=B\)，该 Hom 集映射的满射性给出 \(g:B\rightarrow A\) 使 \(fg=1_B\)。再取 \(X=A\)，由于 \(f(gf)=f=f1_A\)，相应 Hom 集映射的单射性给出 \(gf=1_A\)。所以 \(f\) 是同构。
\end{proof}

这里必须保留“对每个 \(X\)”或另证明所选对象足以检验。在群范畴中，从平凡群到任意群都只有一个同态，仅用这个对象不能区分群。对左 \(R\) 模则可以用 \(R\) 检验同构，因为 \(\operatorname{Hom}_R(R,M)\cong U(M)\)，而双射的模同态具有线性的逆。这一特殊结论依赖具体范畴，不能当作所有范畴都有的底集合判据。

\begin{example}[交换群中的自然一元运算]\label{b4:exm:yoneda-abelian-operations}
设 \(U:\mathbf{Ab}\rightarrow\mathbf{Set}\) 为忘却函子。对每个交换群 \(A\)，给一个集合映射 \(t_A:U(A)\rightarrow U(A)\)，并要求对每个群同态 \(f:A\rightarrow B\) 有 \(f(t_A(a))=t_B(f(a))\)。则存在唯一整数 \(n\)，使所有 \(A\) 和 \(a\in A\) 都满足 \(t_A(a)=na\)。

\ProseParagraphBreak
事实上，取 \(n=t_{\mathbb Z}(1)\)。对唯一满足 \(f_a(1)=a\) 的同态 \(f_a:\mathbb Z\rightarrow A\)，自然性给出 \(t_A(a)=f_a(n)=na\)。反过来，乘固定整数与所有群同态交换。这里从未事先要求 \(t_A\) 为群同态；这一性质被对所有群的自然性迫出。用 Yoneda 语言，这是 \(U\cong h^{\mathbb Z}\) 与 \(\operatorname{Nat}(h^{\mathbb Z},U)\cong U(\mathbb Z)\) 的具体计算。
\end{example}

若只研究一个群 \(A\)，它的任意集合自映射都可以被指定；把这些映射要求为对全部群同态自然，选择就缩减为一个整数。Yoneda 引理的用途正在于把这种同时涉及许多对象的相容条件，转化成一个明确的参数集合。

% !TeX root = ../../Book4.tex
% 纲要标题：### 3.3 伴随函子
% 纲要位置：计划纲要.md，第 3171 行。

\section{伴随函子}
\label{b4:sec:0303}

自由构造与忘却构造之间的泛性质同时涉及两个变化的对象。伴随把这类对应组织成自然的 Hom 双射，而单位、余单位记录双射在恒等态射处的值。

% 正文按纲要写入；各节的基础结果在本章证明。

% 纲要内容（仅作写作计划，尚非教材正文）：
%
% * Hom 集伴随及单位、余单位
% * 三角恒等式
% * 自由—忘却及 Hom–张量伴随
%

\subsection{Hom 集伴随及单位、余单位}
\label{b4:subsec:0303:01}

自由模的可表性针对一个固定集合 \(X\)。当 \(X\) 也变化时，表示对象 \(R^{(X)}\) 本身组成函子，泛性质就同时联系两个范畴。伴随正是这种双变量的自然对应；它不要求两个范畴等价，也不要求某一方忘掉的全部信息都能恢复。

\begin{definition}\label{b4:def:adjunction}
设 \(\mathcal C,\mathcal D\) 为局部小范畴，\(F:\mathcal C\rightarrow\mathcal D\)、\(G:\mathcal D\rightarrow\mathcal C\) 为函子。若存在关于 \(X\in\mathcal C\)、\(Y\in\mathcal D\) 自然的双射
\[
\Phi_{X,Y}:\operatorname{Hom}_{\mathcal D}(FX,Y)
\xrightarrow{\sim}\operatorname{Hom}_{\mathcal C}(X,GY),
\]
则称 \((F,G,\Phi)\) 为一个\term[bansui]{伴随}[adjunction]，称 \(F\) 为 \(G\) 的\term[zuobansui]{左伴随}[left adjoint]，\(G\) 为 \(F\) 的\term[youbansui]{右伴随}[right adjoint]，记 \(F\dashv G\)。这里两边的 Hom 都在源变量 \(X\) 上反变，在目标变量 \(Y\) 上协变：对 \(a:X'\rightarrow X\)、\(b:Y\rightarrow Y'\) 及 \(h:FX\rightarrow Y\)，分别要求
\[
\begin{aligned}
\Phi_{X',Y}\bigl(h\circ F(a)\bigr)
&=\Phi_{X,Y}(h)\circ a,\\
\Phi_{X,Y'}\bigl(b\circ h\bigr)
&=G(b)\circ\Phi_{X,Y}(h).
\end{aligned}
\]
第一式比较沿 \(a\) 的预合成，第二式比较沿 \(b\) 的后合成。两式合并，等价于要求
\begin{equation}\label{b4:eq:adjunction-naturality}
\Phi_{X',Y'}\bigl(b\circ h\circ F(a)\bigr)
=G(b)\circ\Phi_{X,Y}(h)\circ a.
\end{equation}
\end{definition}
\symidx{\(F\dashv G\)：\(F\) 为 \(G\) 的左伴随}

\[
\begin{tikzcd}[column sep=large,row sep=4.5em]
\operatorname{Hom}_{\mathcal D}(FX,Y)
\arrow[r,"\Phi_{X,Y}","\sim"']
\arrow[d,"b\circ(-)\circ F(a)" description]&
\operatorname{Hom}_{\mathcal C}(X,GY)
\arrow[d,"G(b)\circ(-)\circ a" description]\\
\operatorname{Hom}_{\mathcal D}(FX',Y')
\arrow[r,"\Phi_{X',Y'}","\sim"']&
\operatorname{Hom}_{\mathcal C}(X',GY').
\end{tikzcd}
\]

固定 \(X\) 后，\(FX\) 表示协变函子 \(Y\mapsto\operatorname{Hom}_{\mathcal C}(X,GY)\)；固定 \(Y\) 后，\(GY\) 表示反变函子 \(X\mapsto\operatorname{Hom}_{\mathcal D}(FX,Y)\)。这分别是\secref{b4:sec:0301}中的协变表示与反变表示。与只找一个表示对象相比，伴随还要求这些表示在另一个变量变化时自然相容；函子 \(F,G\) 不仅给出表示对象，也把表示对象之间的态射组织起来。因此不能只对每对对象随意选取集合双射。

\ProseParagraphBreak
泛元素判据说明，表示对象上的恒等态射所对应的元素决定全部表示双射。把这个思想用于上述两族表示，两个恒等态射便给出映射：
\begin{equation}\label{b4:eq:adjunction-unit-counit-def}
\eta_X\coloneqq\Phi_{X,FX}(1_{FX}):X\longrightarrow GFX,
\qquad
\varepsilon_Y\coloneqq\Phi^{-1}_{GY,Y}(1_{GY}):FGY\longrightarrow Y.
\end{equation}
它们分别称为伴随的\term[bansui-danwei]{单位}[unit]与\term[bansui-yudanwei]{余单位}[counit]。\(\eta_X\) 正是固定 \(X\) 后由 \(FX\) 给出的协变表示的泛元素，\(\varepsilon_Y\) 则是固定 \(Y\) 后由 \(GY\) 给出的反变表示的泛元素。单位将一个对象送入“先作 \(F\) 再作 \(G\)”得到的对象；余单位则从“先作 \(G\) 再作 \(F\)”的对象回到原对象。
\symidx{\(\eta,\ \varepsilon\)：伴随的单位与余单位}

\begin{proposition}\label{b4:prop:adjunction-transpose-formulas}
设 \(\mathcal C,\mathcal D\) 为局部小范畴，\(F:\mathcal C\rightarrow\mathcal D\)、\(G:\mathcal D\rightarrow\mathcal C\) 为伴随函子 \(F\dashv G\)，其自然双射为 \(\Phi_{X,Y}:\operatorname{Hom}_{\mathcal D}(FX,Y)\xrightarrow{\sim}\operatorname{Hom}_{\mathcal C}(X,GY)\)。定义 \(\eta_X=\Phi_{X,FX}(1_{FX})\)、\(\varepsilon_Y=\Phi^{-1}_{GY,Y}(1_{GY})\)。这些映射组成自然变换 \(\eta:1_{\mathcal C}\Rightarrow GF\) 与 \(\varepsilon:FG\Rightarrow1_{\mathcal D}\)，且
\begin{equation}\label{b4:eq:adjunction-transpose-formulas}
\Phi_{X,Y}(h)=G(h)\circ\eta_X,
\qquad
\Phi^{-1}_{X,Y}(k)=\varepsilon_Y\circ F(k)
\end{equation}
对所有 \(h:FX\rightarrow Y\)、\(k:X\rightarrow GY\) 成立。
\end{proposition}
\begin{proof}
对 \(a:X\rightarrow X'\)，要证明单位的自然性，就要比较 \(GF(a)\circ\eta_X\) 与 \(\eta_{X'}\circ a\)。前者来自双射在第二变量上的自然性，后者来自第一变量上的自然性。为使两种自然性作用于同一个态射，把 \(F(a):FX\rightarrow FX'\) 分别写作 \(F(a)\circ1_{FX}\) 与 \(1_{FX'}\circ F(a)\)，得到
\[
GF(a)\circ\eta_X
=\Phi_{X,FX'}\bigl(F(a)\bigr)
=\eta_{X'}\circ a.
\]
所以 \(\eta\) 自然。因每个 \(\Phi_{X,Y}\) 为双射，将\cref{b4:eq:adjunction-naturality}反向改写可知，逆双射也满足两变量自然性：对 \(a:X'\rightarrow X\)、\(b:Y\rightarrow Y'\) 及 \(k:X\rightarrow GY\)，有
\[
\Phi^{-1}_{X',Y'}\bigl(G(b)\circ k\circ a\bigr)
=b\circ\Phi^{-1}_{X,Y}(k)\circ F(a).
\]
对 \(b:Y\rightarrow Y'\)，余单位的自然性要求比较 \(b\circ\varepsilon_Y\) 与 \(\varepsilon_{Y'}\circ FG(b)\)。同样，在逆双射下分别把 \(G(b)\) 写成 \(G(b)\circ1_{GY}\) 与 \(1_{GY'}\circ G(b)\)，让第二变量与第一变量的自然性分别作用，得到
\[
b\circ\varepsilon_Y
=\Phi^{-1}_{GY,Y'}\bigl(G(b)\bigr)
=\varepsilon_{Y'}\circ FG(b).
\]
这证明 \(\varepsilon\) 自然。最后，把 \(h\) 写成 \(h\circ1_{FX}\)，在第二变量使用自然性，得到第一条转置公式；把 \(k\) 写成 \(1_{GY}\circ k\)，对逆双射在第一变量使用自然性，得到第二条公式。
\end{proof}

因此，对每个 \(\mathcal C\) 中的态射 \(k:X\rightarrow GY\)，存在唯一的 \(\mathcal D\) 中态射 \(h:FX\rightarrow Y\)，使 \(k=G(h)\eta_X\)。单位 \(\eta_X\) 因而为 \(X\) 到所有 \(GY\) 的映射提供一个统一的分解：先到 \(GFX\)，再用某个唯一确定的 \(h\) 所诱导的 \(G(h)\)。对偶地，每个 \(\mathcal D\) 中的态射 \(h:FX\rightarrow Y\) 唯一写成 \(h=\varepsilon_YF(k)\)，其中 \(k:X\rightarrow GY\)；来自所有 \(FX\) 的映射都经 \(FGY\) 后用同一个 \(\varepsilon_Y\) 到达 \(Y\)。这分别是单位与余单位本身的泛性质：

\[
\begin{tikzcd}[column sep=large,row sep=large]
X\arrow[r,"\eta_X"]\arrow[dr,"k"']&GFX\arrow[d,dashed,"G(h)"]\\
&GY
\end{tikzcd}
\qquad
\begin{gathered}
h\in\operatorname{Hom}_{\mathcal D}(FX,Y),\\
\text{满足左图交换的}\;h\;\text{唯一。}
\end{gathered}
\]
\[
\begin{tikzcd}[column sep=large,row sep=large]
FX\arrow[r,dashed,"F(k)"]\arrow[dr,"h"']&FGY\arrow[d,"\varepsilon_Y"]\\
&Y.
\end{tikzcd}
\qquad
\begin{gathered}
k\in\operatorname{Hom}_{\mathcal C}(X,GY),\\
\text{满足左图交换的}\;k\;\text{唯一。}
\end{gathered}
\]

第二卷第6.3节中“指定纯张量的像”和“固定一个变量”，正是这两个公式在模范畴中的表现。

\subsection{三角恒等式}
\label{b4:subsec:0303:02}

从单位和余单位能否恢复整个伴随？两条转置公式已经给出候选映射；所需条件就是保证它们互逆。称下列两式为\term[sanjiao-hengdengshi]{三角恒等式}[triangle identities]：
\begin{equation}\label{b4:eq:triangle-identities}
\varepsilon_{FX}\circ F(\eta_X)=1_{FX},
\qquad
G(\varepsilon_Y)\circ\eta_{GY}=1_{GY}.
\end{equation}
第一式保证从 \(h:FX\rightarrow Y\) 出发，转置为 \(X\rightarrow GY\) 再逆转置时回到 \(h\)；第二式保证从 \(k:X\rightarrow GY\) 出发，逆转置为 \(FX\rightarrow Y\) 再转置时回到 \(k\)。它们分别把 \(FX\) 或 \(GY\) 经三重函子复合再返回原对象的复合消为恒等态射。

\[
\begin{tikzcd}[column sep=large,row sep=large]
FX\arrow[r,"F(\eta_X)"]\arrow[dr,"1_{FX}"']&
FGFX\arrow[d,"\varepsilon_{FX}"]\\
&FX
\end{tikzcd}
\]
\[
\begin{tikzcd}[column sep=large,row sep=large]
GY\arrow[r,"\eta_{GY}"]\arrow[dr,"1_{GY}"']&
GFGY\arrow[d,"G(\varepsilon_Y)"]\\
&GY.
\end{tikzcd}
\]

\begin{theorem}[伴随的单位与余单位刻画]\label{b4:thm:adjunction-unit-counit}
设 \(\mathcal C,\mathcal D\) 为局部小范畴，\(F:\mathcal C\rightarrow\mathcal D\)、\(G:\mathcal D\rightarrow\mathcal C\) 为函子。关于两变量自然的伴随双射
\[
\operatorname{Hom}_{\mathcal D}(FX,Y)
\cong\operatorname{Hom}_{\mathcal C}(X,GY)
\]
与满足三角恒等式
\[
\varepsilon_{FX}F(\eta_X)=1_{FX}\quad(X\in\mathcal C),\qquad
G(\varepsilon_Y)\eta_{GY}=1_{GY}\quad(Y\in\mathcal D)
\]
的自然变换对
\(\eta:1_{\mathcal C}\Rightarrow GF\)、\(\varepsilon:FG\Rightarrow1_{\mathcal D}\)
一一对应。由双射 \(\Phi\) 得到变换的公式为 \(\eta_X=\Phi_{X,FX}(1_{FX})\)、\(\varepsilon_Y=\Phi^{-1}_{GY,Y}(1_{GY})\)；反向得到双射及其逆的公式为 \(\Phi_{X,Y}(h)=G(h)\eta_X\)、\(\Phi^{-1}_{X,Y}(k)=\varepsilon_YF(k)\)，其中 \(h:FX\rightarrow Y\)、\(k:X\rightarrow GY\)。
\end{theorem}
\begin{proof}
先给定伴随双射。单位、余单位的自然性已经由\cref{b4:prop:adjunction-transpose-formulas}证明。将第一条转置公式所得的 \(\eta_X=\Phi(1_{FX})\) 再代入逆公式，得到
\[
1_{FX}=\Phi^{-1}(\eta_X)=\varepsilon_{FX}F(\eta_X).
\]
将 \(\varepsilon_Y=\Phi^{-1}(1_{GY})\) 代回正向公式，得到
\(1_{GY}=G(\varepsilon_Y)\eta_{GY}\)。所以三角恒等式成立。

反过来，设 \(\eta,\varepsilon\) 自然且满足三角恒等式，定义
\[
\Phi(h)=G(h)\eta_X,\qquad \Psi(k)=\varepsilon_YF(k).
\]
对于 \(k:X\rightarrow GY\)，\(\eta\) 关于 \(k\) 的自然性给出 \(GF(k)\eta_X=\eta_{GY}k\)，把 \(k\) 移到复合的右端；此后中间的 \(G(\varepsilon_Y)\eta_{GY}\) 正好可由第二条三角恒等式消去：
\[
\Phi\Psi(k)
=G(\varepsilon_Y)GF(k)\eta_X
=G(\varepsilon_Y)\eta_{GY}k=k.
\]
对于 \(h:FX\rightarrow Y\)，\(\varepsilon\) 关于 \(h\) 的自然性给出 \(\varepsilon_YFG(h)=h\varepsilon_{FX}\)，把 \(h\) 移到复合的左端；此后 \(\varepsilon_{FX}F(\eta_X)\) 由第一条三角恒等式消去：
\[
\Psi\Phi(h)
=\varepsilon_YFG(h)F(\eta_X)
=h\varepsilon_{FX}F(\eta_X)=h.
\]
因此 \(\Phi,\Psi\) 互逆。再对 \(a:X'\rightarrow X\)、\(b:Y\rightarrow Y'\) 计算
\[
\begin{aligned}
\Phi\bigl(bhF(a)\bigr)
&=G(b)G(h)GF(a)\eta_{X'}\\
&=G(b)G(h)\eta_Xa
=G(b)\Phi(h)a,
\end{aligned}
\]
证明双射在两变量上自然。由新双射在恒等态射上的取值，重新得到的单位与余单位仍为 \(\eta,\varepsilon\)。另一方向则由已证转置公式恢复原双射，故两种构造互逆。
\end{proof}

可对照 Weibel 的伴随定义及单位、余单位定理\cite[A.6.1; Theorem A.6.2 and Exercise A.6.2]{Weibel1994HomologicalAlgebra}。

\ProseParagraphBreak
三角恒等式不意味着 \(\eta_X\) 与 \(\varepsilon_Y\) 互逆：它们通常连定义域、陪域都不同。例如自由交换群伴随中，一个集合被送到它的形式整数线性组合；回到具体群时，余单位把这些形式和求值。这个过程通常带来许多关系，远非对象不变的同构。

\subsection{自由—忘却及 Hom–张量伴随}
\label{b4:subsec:0303:03}

设 \(R\) 为含幺环。自由左模函子 \(F_R:\mathbf{Set}\rightarrow R\text{-}\mathbf{Mod}\) 将 \(X\) 送到 \(R^{(X)}\)，并把函数 \(a:X\rightarrow X'\) 送到唯一满足 \(e_x\mapsto e_{a(x)}\) 的模同态。恒等与复合律由各基向量的像确定。\secref{b4:sec:0301}给出的双射不仅在模变量上自然，在集合变量上也自然：对 \(f:R^{(X')}\rightarrow M\)，前复合 \(F_R(a)\) 对应函数 \(x\mapsto f(e_{a(x)})\)。所以
\[
F_R\dashv U.
\]
单位为 \(\eta_X(x)=e_x\)，余单位为
\[
\varepsilon_M:R^{(U(M))}\longrightarrow M,
\qquad\sum_mr_me_m\longmapsto\sum_mr_mm.
\]
这里 \(e_m\) 是以 \(m\in M\) 为标签的自由模基向量，并非把 \(m\) 本身当作自由模中的元素。对于 \(x\in X\)，\(e_x\) 属于 \(F_RX\)，而 \(e_{e_x}\) 属于 \(F_RU(F_RX)\)：后者是以自由模元素 \(e_x\) 为标签的新基向量。第一条三角恒等式在生成元上的作用为
\[
e_x\longmapsto e_{e_x}\longmapsto e_x.
\]
第二条在 \(U(M)\) 上的作用为 \(m\mapsto e_m\mapsto m\)。两式分别是在适当的一层插入形式基向量后再求值。

\ProseParagraphBreak
另一组伴随来自固定双模。它是第二卷第6.3节 Hom–张量伴随定理的原有左模版本；本卷把它放入已建立的一般定义，不重新构造张量积。

\begin{proposition}\label{b4:prop:hom-tensor-adjunction}
设 \(R,S\) 为含幺环，\({}_SP_R\) 为幺性左 \(S\) 右 \(R\) 双模。以下模范畴的对象均为幺性左模。函子
\[
F=P\otimes_R-:R\text{-}\mathbf{Mod}\longrightarrow S\text{-}\mathbf{Mod},
\qquad
G=\operatorname{Hom}_S(P,-):S\text{-}\mathbf{Mod}\longrightarrow R\text{-}\mathbf{Mod}
\]
满足 \(F\dashv G\)，其中 \(G(N)\) 的左 \(R\) 作用为 \((rf)(p)=f(pr)\)。对任意幺性左 \(R\) 模 \(M\) 与幺性左 \(S\) 模 \(N\)，其单位与余单位为
\[
\eta_M(m)(p)=p\otimes m,
\qquad
\varepsilon_N(p\otimes f)=f(p).
\]
\end{proposition}
\begin{proof}
第二卷第6.3节的定理对任意左 \(R\) 模 \(M\) 和左 \(S\) 模 \(N\) 给出自然双射
\[
h:P\otimes_RM\rightarrow N
\quad\longleftrightarrow\quad
\bigl(m\mapsto(p\mapsto h(p\otimes m))\bigr).
\]
反向地，一个左 \(R\) 模同态 \(k:M\rightarrow\operatorname{Hom}_S(P,N)\) 对应
\[
\widetilde k:P\otimes_RM\longrightarrow N,
\qquad p\otimes m\longmapsto k(m)(p).
\]
这个公式的平衡性由
\(k(rm)(p)=(rk(m))(p)=k(m)(pr)\) 保证，左 \(S\) 线性由 \(k(m)\in\operatorname{Hom}_S(P,N)\) 保证。两次转置在 \(m,p\) 上恢复原来的值，故互逆。第二卷还证明了两变量自然性，且所用环和双模假设与本命题完全相同，故满足\cref{b4:def:adjunction}。将两个恒等映射代入便得到所列单位与余单位。固定 \(m\) 时，\(p\mapsto p\otimes m\) 为加性映射，且
\[
\eta_M(m)(sp)=sp\otimes m=s\cdot(p\otimes m),
\]
所以 \(\eta_M(m)\) 确实属于 \(\operatorname{Hom}_S(P,P\otimes_RM)\)。再把 \(m\) 当作变量，张量积的加性给出 \(\eta_M\) 的加性，而左 \(R\) 作用满足
\[
\eta_M(rm)(p)=p\otimes rm=pr\otimes m
=\bigl(r\eta_M(m)\bigr)(p),
\]
所以单位左 \(R\) 线性；余单位的平衡性为 \(f(pr)=(rf)(p)\)，其左 \(S\) 线性为 \(f(sp)=s\cdot f(p)\)。第一条三角恒等式在纯张量上将 \(p\otimes m\) 送到 \(\eta_M(m)(p)=p\otimes m\)；第二条将 \(f\) 送到函数 \(p\mapsto\varepsilon_N(p\otimes f)=f(p)\)。因此这两个具体映射也直接满足一般定理的条件。
\end{proof}

本命题只规定 \(N\) 是左 \(S\) 模，没有赋予它左 \(R\) 作用，因此 \(r\cdot f(p)\) 通常根本没有定义。\((rf)(p)=f(pr)\) 使用的是 \(P\) 的右 \(R\) 作用；这个作用经输入端传到 Hom 上，才使其成为伴随所需的左 \(R\) 模。

\begin{example}[交换化与包含函子]\label{b4:exm:abelianization-adjunction}
对群 \(H\)，记 \(H^{\mathrm{ab}}=H/[H,H]\)，其中 \([H,H]\) 是由交换子生成的正规子群。每个到交换群 \(A\) 的群同态都消去交换子，因而唯一通过商群分解，给出
\[
\operatorname{Hom}_{\mathbf{Ab}}(H^{\mathrm{ab}},A)
\cong\operatorname{Hom}_{\mathbf{Grp}}(H,A).
\]
前复合群同态或后复合交换群同态都与商映射分解相容，因而得到伴随
\[
(-)^{\mathrm{ab}}:\mathbf{Grp}\longrightarrow\mathbf{Ab}
\quad\dashv\quad
I:\mathbf{Ab}\hookrightarrow\mathbf{Grp},
\]
其中 \(I\) 是包含函子。单位为商映射 \(H\rightarrow H^{\mathrm{ab}}\)；因 \(A\) 已经交换，\([A,A]=\{1\}\)，余单位 \(A^{\mathrm{ab}}\xrightarrow{\sim}A\) 是同构。非交换群的单位不是单射，而余单位总是同构。这也说明伴随不要求单位与余单位具有相同的可逆性。
\end{example}

% !TeX root = ../../Book4.tex
% 纲要标题：### 3.4 伴随的后果
% 纲要位置：计划纲要.md，第 3177 行。

\section{伴随的后果}
\label{b4:sec:0304}

伴随双射可以直接改写极限或余极限的泛性质，从而确定哪些构造被保留。单位与余单位的复合还产生单子结构，说明连续施加伴随函子时出现的代数规律。

% 正文按纲要写入；各节的基础结果在本章证明。

% 纲要内容（仅作写作计划，尚非教材正文）：
%
% * 左伴随保持余极限
% * 右伴随保持极限
% * 一般伴随函子定理的极限保持、解集和选择条件，完整证明置于附录
% * 单子与余单子初步，完整单子性理论置于附录
% ---
%

\subsection{左伴随与余极限}
\label{b4:subsec:0304:01}

余极限的泛性质通过从所得对象出发的映射表述，伴随双射恰好能改写这种映射。下面的结论因此不依赖余极限在某个范畴中的具体构造。

\begin{theorem}\label{b4:thm:adjoints-preserve-colimits}
设 \(\mathcal C,\mathcal D\) 为局部小范畴，\(F:\mathcal C\rightarrow\mathcal D\) 有右伴随 \(G:\mathcal D\rightarrow\mathcal C\)。若小范畴 \(J\) 上的图表 \(K:J\rightarrow\mathcal C\) 有余极限 \((L,(\iota_j)_{j\in J})\)，则 \((FL,(F\iota_j)_{j\in J})\) 是 \(FK\) 的余极限。因此左伴随保持所有存在的小余极限，包括空图表的余极限。
\end{theorem}
\begin{proof}
记伴随在两变量自然的双射为
\[
\Phi_{X,Y}:\operatorname{Hom}_{\mathcal D}(FX,Y)
\xrightarrow{\sim}\operatorname{Hom}_{\mathcal C}(X,GY),
\qquad X\in\mathcal C,\ Y\in\mathcal D.
\]
函子保持复合，故 \((F\iota_j)\) 是余锥。要验证它的泛性质，可以依次改写从 \(FL\) 出发的 Hom 集：
\[
\begin{aligned}
\operatorname{Hom}_{\mathcal D}(FL,Y)
&\cong\operatorname{Hom}_{\mathcal C}(L,GY)\\
&\cong\{\,K\;\text{到}\;GY\;\text{的余锥}\,\}\\
&\cong\{\,FK\;\text{到}\;Y\;\text{的余锥}\,\}.
\end{aligned}
\]
第一步是伴随双射，第二步是 \(L\) 的余极限泛性质，第三步逐分量使用伴随双射；自然性保证余锥的相容条件也随之对应。这些双射关于 \(Y\) 自然。具体地，取余锥 \(q_j:FK(j)\rightarrow Y\)，令
\[
s_j=\Phi_{K(j),Y}(q_j):K(j)\longrightarrow GY.
\]
对 \(a:i\rightarrow j\)，余锥相容性为 \(q_jFK(a)=q_i\)。伴随在第一变量的自然性给出 \(s_jK(a)=s_i\)，所以 \((s_j)\) 是 \(K\) 的余锥。由 \(L\) 的泛性质，存在唯一 \(s:L\rightarrow GY\) 满足 \(s\iota_j=s_j\)。令 \(q=\Phi^{-1}_{L,Y}(s):FL\rightarrow Y\)。再次利用自然性，
\[
\Phi_{K(j),Y}\bigl(qF(\iota_j)\bigr)
=\Phi_{L,Y}(q)\iota_j=s\iota_j=s_j=\Phi_{K(j),Y}(q_j),
\]
双射性给出 \(qF(\iota_j)=q_j\)。若 \(q'\) 也满足这些等式，则 \(\Phi(q')\) 也经过每个 \(\iota_j\) 给出 \(s_j\)，由余极限的唯一性有 \(\Phi(q')=s\)，从而 \(q'=q\)。证明没有要求 \(J\) 非空；在空图表情形，它就是始对象的泛性质。
\end{proof}

例如 \({}_SP_R\) 固定时，\cref{b4:prop:hom-tensor-adjunction}表明 \(P\otimes_R-\) 保持任意直和、余等化子及滤过余极限。这并非说所有相关对象的底集合都按同一公式构造，而是说它们具有同一个余极限泛性质。具体的比较同构由结构映射唯一确定，例如
\[
\bigoplus_i(P\otimes_RM_i)\xrightarrow{\sim}P\otimes_R\left(\bigoplus_iM_i\right)
\]
在第 \(i\) 个直和项上把 \(p\otimes m\) 送到 \(p\otimes\iota_i(m)\)。

\ProseParagraphBreak
左伴随的条件本身不保证保持积。\cref{b4:exm:exact-not-products}已经用共同分母证明，\(\mathbb Q\otimes_{\mathbb Z}-:\mathbf{Ab}\rightarrow\mathbb Q\text{-}\mathbf{Mod}\) 不保持可数个 \(\mathbb Z\) 的积。现在将\cref{b4:prop:hom-tensor-adjunction}用于 \({}_{\mathbb Q}\mathbb Q_{\mathbb Z}\)，又知道它是左伴随，因此这个例子也检验了保持定理的方向：它保持余极限，却仍然可能不保持无限积。此处不能从无限积失败推断有限积也失败。

\subsection{右伴随与极限}
\label{b4:subsec:0304:02}

\begin{theorem}\label{b4:thm:adjoints-preserve-limits}
设 \(\mathcal C,\mathcal D\) 为局部小范畴，\(F:\mathcal C\rightarrow\mathcal D\) 左伴随于 \(G:\mathcal D\rightarrow\mathcal C\)。若小图表 \(K:J\rightarrow\mathcal D\) 有极限 \((L,(p_j)_{j\in J})\)，则 \((GL,(Gp_j)_{j\in J})\) 是 \(GK\) 的极限。因此右伴随保持所有存在的小极限，包括空图表的极限。
\end{theorem}
\begin{proof}
记伴随在两变量自然的双射为
\[
\Phi_{X,Y}:\operatorname{Hom}_{\mathcal D}(FX,Y)
\xrightarrow{\sim}\operatorname{Hom}_{\mathcal C}(X,GY),
\qquad X\in\mathcal C,\ Y\in\mathcal D.
\]
这次改写的是指向极限的 Hom 集：
\[
\begin{aligned}
\operatorname{Hom}_{\mathcal C}(X,GL)
&\cong\operatorname{Hom}_{\mathcal D}(FX,L)\\
&\cong\{\,FX\;\text{到}\;K\;\text{的锥}\,\}\\
&\cong\{\,X\;\text{到}\;GK\;\text{的锥}\,\}.
\end{aligned}
\]
中间一步使用 \(L\) 的极限泛性质，两侧使用伴随双射，所有对应关于 \(X\) 自然。具体地，给定锥 \(s_j:X\rightarrow GK(j)\)，用逆伴随双射得到 \(q_j:FX\rightarrow K(j)\)。对 \(a:i\rightarrow j\)，关系 \(GK(a)s_i=s_j\) 通过第二变量自然性等价于 \(K(a)q_i=q_j\)。因此 \((q_j)\) 为锥，唯一确定 \(q:FX\rightarrow L\)，使 \(p_jq=q_j\)。令 \(s=\Phi_{X,L}(q)\)，第二变量自然性给出
\[
Gp_j\circ s=\Phi_{X,K(j)}(p_jq)=s_j.
\]
任何具有这些复合的 \(s'\)，在逆双射下都给出具有投影 \(q_j\) 的态射 \(q'\)，由 \(L\) 的唯一性有 \(q'=q\)，进而 \(s'=s\)。空图表情形同样成立，给出终对象的保持。
\end{proof}

固定源的 Hom 函子保持极限已在\cref{b4:prop:representables-preserve-limits}直接证明：对于有极限的图表 \(K:J\rightarrow\mathcal C\) 及 \(A\in\mathcal C\)，一个态射 \(A\rightarrow\lim_JK\) 恰好对应从 \(A\) 到各 \(K(j)\) 的相容态射族，所以
\[
\operatorname{Hom}_{\mathcal C}(A,\lim_JK)
\cong\lim_{j\in J}\operatorname{Hom}_{\mathcal C}(A,K(j)).
\]
若一个函子与固定源的 Hom 函子自然同构，相应的锥经分量同构转换，仍具有相同的唯一分解性质，所以任意协变可表函子保持极限。反变可表函子则把原范畴的余极限变成集合的极限；必须连同相反范畴一起读这个断言，不能遗漏方向。

\ProseParagraphBreak
由 Hom–张量伴随，\(\operatorname{Hom}_S(P,-)\) 保持积与等化子。若给定左 \(S\) 模同态 \(u:N\rightarrow N'\)，它的核是 \(u\) 与零映射的等化子，故
\[
\operatorname{Hom}_S(P,\ker u)
\cong\ker\bigl(\operatorname{Hom}_S(P,u)\bigr).
\]
右边由那些满足 \(u\circ f=0\) 的 \(f:P\rightarrow N\) 组成。一般情形下，\(\operatorname{Hom}_S(P,-)\) 未必保持满同态；例如 \(\mathbb Z\rightarrow\mathbb Z/n\mathbb Z\) 在 \(n>1\) 时满射，而
\[
\operatorname{Hom}_{\mathbb Z}(\mathbb Z/n\mathbb Z,\mathbb Z)=0,
\qquad
\operatorname{Hom}_{\mathbb Z}(\mathbb Z/n\mathbb Z,\mathbb Z/n\mathbb Z)
\cong\mathbb Z/n\mathbb Z.
\]
因此诱导映射不是满射。保持核与有限积使这个 Hom 函子具有\secref{b4:sec:0204}中的左正合性；上例的满射端失败，说明它一般不右正合。下一章会把这些正合性质放到一般交换范畴中讨论。

\ProseParagraphBreak
保持极限是一项必要条件，并非任意情况下都足以推出存在左伴随。证明存在伴随还要为每个对象构造相应的表示对象，并控制候选对象的大小。设 \(\mathcal C,\mathcal D\) 为局部小范畴，\(\mathcal D\) 有全部小极限，\(U:\mathcal D\to\mathcal C\) 为函子。解集条件要求：对每个 \(X\in\mathcal C\)，存在集合大小的一族箭头 \(u_i:X\to UY_i\)，使每个 \(u:X\to UY\) 都能写成 \(u=U(v)u_i\)，其中 \(v:Y_i\to Y\)。它不要求这批候选已经具有唯一分解性质，而是从可能遍历真类的候选中选出集合大小的一族，足以产生全部这些箭头。再假定满足解集条件时，所需的解集与小极限可以按对象类同时选择。\Cref{b4:thm:A-general-adjoint-functor}将在这些条件下证明：\(U\) 有左伴随，当且仅当它保持全部小极限并满足解集条件。

\subsection{单子与余单子初步}
\label{b4:subsec:0304:03}

对伴随 \(F\dashv G\)，复合 \(T=GF\) 是 \(\mathcal C\) 上的自函子；\(T^2\) 表示函子复合 \(T\circ T=GFGF\)。单位给出 \(X\rightarrow TX\)。余单位在 \(FX\) 处的分量 \(\varepsilon_{FX}:FGFX\rightarrow FX\) 将中间的 \(FG\) 归并，施加 \(G\) 后得到
\[
\mu_X=G(\varepsilon_{FX}):T^2X=GFGFX\longrightarrow GFX=TX.
\]
在自由左 \(R\) 模与底集合函子的伴随中，\(TX\) 是 \(X\) 的形式有限线性组合，而 \(T^2X\) 是“形式线性组合的形式线性组合”。\(\mu_X\) 展开括号并合并同类项。这提示我们暂时只保留自函子、插入元素和归并运算。

\begin{definition}\label{b4:def:monad}
设 \(\mathcal C\) 为范畴，\(T:\mathcal C\rightarrow\mathcal C\) 为自函子，\(\eta:1_{\mathcal C}\Rightarrow T\)、\(\mu:T^2\Rightarrow T\) 为自然变换。若对每个 \(X\in\mathcal C\) 都有
\[
\mu_X\circ T(\mu_X)=\mu_X\circ\mu_{TX},
\qquad
\mu_X\circ T(\eta_X)=1_{TX}=\mu_X\circ\eta_{TX},
\]
则称 \((T,\eta,\mu)\) 为 \(\mathcal C\) 上的\term[danzi]{单子}[monad]。第一式称结合律，后两式称单位律。
\end{definition}
\symidx{\(\mu:T^2\Rightarrow T\)：单子的乘法自然变换}

\[
\begin{tikzcd}[column sep=large,row sep=large]
T^3X\arrow[r,"T(\mu_X)"]\arrow[d,"\mu_{TX}"']&
T^2X\arrow[d,"\mu_X"]\\
T^2X\arrow[r,"\mu_X"']&TX
\end{tikzcd}
\]
\[
\begin{tikzcd}[column sep=large,row sep=large]
TX\arrow[r,shift left=0.7ex,"T(\eta_X)"]
\arrow[r,shift right=0.7ex,"\eta_{TX}"']
\arrow[dr,"1_{TX}"']&T^2X\arrow[d,"\mu_X"]\\
&TX.
\end{tikzcd}
\]

结合律的方块比较从三层复合到一层复合的两种归并次序。单位律的图则有两个不同的插入位置：\(T(\eta_X)\) 在已经施加的 \(T\) 内部插入单位，\(\eta_{TX}\) 对整个 \(TX\) 插入单位。二者有相同的源与目标，却未必是同一态射；单位律要求它们分别与 \(\mu_X\) 复合后都还原 \(TX\)。

\begin{proposition}\label{b4:prop:adjunction-monad}
设 \(F:\mathcal C\rightarrow\mathcal D\)、\(G:\mathcal D\rightarrow\mathcal C\) 为局部小范畴之间的伴随，其单位、余单位为 \(\eta,\varepsilon\)。则
\[
T=GF,\qquad \mu_X=G(\varepsilon_{FX})
\]
与原有单位 \(\eta\) 组成 \(\mathcal C\) 上的单子。
\end{proposition}
\begin{proof}
对 \(f:X\rightarrow X'\)，要证明 \(T(f)\mu_X=\mu_{X'}T^2(f)\)。将 \(F(f):FX\rightarrow FX'\) 代入余单位的自然性，得到
\[
F(f)\circ\varepsilon_{FX}
=\varepsilon_{FX'}\circ FG(F(f)).
\]
施加 \(G\)，左边为 \(T(f)\mu_X\)，右边为 \(\mu_{X'}T^2(f)\)，故 \(\mu\) 自然。两条单位律分别为
\[
\begin{aligned}
\mu_XT(\eta_X)
&=G\bigl(\varepsilon_{FX}F(\eta_X)\bigr)=1_{GFX},\\
\mu_X\eta_{TX}
&=G(\varepsilon_{FX})\eta_{GFX}=1_{GFX},
\end{aligned}
\]
第一行对 \(\varepsilon_{FX}F(\eta_X)=1_{FX}\) 施加 \(G\)，第二行将另一个三角恒等式中的 \(Y\) 取为 \(FX\)。

结合律也来自余单位的自然性，但这次使用的态射是
\(u=\varepsilon_{FX}:FGFX\rightarrow FX\)。对这两个对象及 \(u\)，自然性给出
\[
u\circ\varepsilon_{FGFX}
=\varepsilon_{FX}\circ FG(u).
\]
施加 \(G\) 后，\(G(u)=\mu_X\)，而
\[
G(\varepsilon_{FGFX})=\mu_{TX},\qquad
G(FG(u))=T(\mu_X).
\]
因此得到 \(\mu_X\mu_{TX}=\mu_XT(\mu_X)\)，正是结合律。
\end{proof}

在自由左 \(R\) 模伴随中，设 \(R\) 为任意含幺环。令 \(e_x\) 表示 \(TX\) 中的形式生成元，而 \(e_{v_i}\) 表示 \(T^2X\) 中由 \(v_i\in TX\) 标记的生成元，则
\[
\mu_X\left(\sum_i r_i e_{v_i}\right)
=\sum_i r_iv_i
=\sum_x\left(\sum_i r_is_{i,x}\right)e_x,
\qquad v_i=\sum_xs_{i,x}e_x.
\]
所有和都有限，标量次序为 \(r_is_{i,x}\)，不能在非交换环中调换。对 \(v=\sum_xs_xe_x\in TX\)，两种插入方式具体为
\[
T(\eta_X)(v)=\sum_xs_xe_{e_x},\qquad
\eta_{TX}(v)=e_v.
\]
前者把每个原生成元再包成一项，后者把整个线性组合当成一个新生成元；\(\mu_X\) 分别作用后都得到 \(v\)。

\ProseParagraphBreak
三层的结合律也可以直接计算。取 \(X=\{x,y,z\}\) 及 \(r,s,t\in R\)，令
\[
v=s e_x+e_y\in TX,\qquad
w=r e_v+e_{e_z}\in T^2X,\qquad
\xi=t e_w\in T^3X.
\]
按两种次序归并，分别得到
\[
\begin{aligned}
\mu_XT(\mu_X)(\xi)
&=\mu_X\bigl(t e_{rs e_x+r e_y+e_z}\bigr)\\
&=t(rs)e_x+(tr)e_y+t e_z,\\
\mu_X\mu_{TX}(\xi)
&=\mu_X\bigl((tr)e_v+t e_{e_z}\bigr)\\
&=((tr)s)e_x+(tr)e_y+t e_z.
\end{aligned}
\]
环的结合律给出 \(t(rs)=(tr)s\)，所以两条路线的结果相同，且整个计算没有交换任何标量的次序。

\begin{definition}\label{b4:def:comonad}
设 \(\mathcal D\) 为范畴，\(Q:\mathcal D\rightarrow\mathcal D\) 为自函子，\(\varepsilon:Q\Rightarrow1_{\mathcal D}\)、\(\delta:Q\Rightarrow Q^2\) 为自然变换。若对每个 \(Y\in\mathcal D\) 都有
\[
Q(\delta_Y)\circ\delta_Y=\delta_{QY}\circ\delta_Y,
\qquad
Q(\varepsilon_Y)\circ\delta_Y=1_{QY}
=\varepsilon_{QY}\circ\delta_Y,
\]
则称 \((Q,\varepsilon,\delta)\) 为 \(\mathcal D\) 上的\term[yudanzi]{余单子}[comonad]。
\end{definition}
\symidx{\(\delta:Q\Rightarrow Q^2\)：余单子的余乘法自然变换}

伴随 \(F\dashv G\) 在另一范畴上给出 \(Q=FG\)，余单位仍为 \(\varepsilon\)。这次用单位插入中间一层：
\[
\delta_Y=F(\eta_{GY}):QY=FGY\longrightarrow FGFGY=Q^2Y.
\]
其自然性由 \(\eta\) 的自然性经 \(F\) 得到。两条余单位律可以直接核验：
\[
\begin{aligned}
\varepsilon_{QY}\circ\delta_Y
&=\varepsilon_{FGY}\circ F(\eta_{GY})=1_{FGY},\\
Q(\varepsilon_Y)\circ\delta_Y
&=F\bigl(G(\varepsilon_Y)\circ\eta_{GY}\bigr)=1_{FGY}.
\end{aligned}
\]
第一行使用 \(GY\) 处的三角恒等式，第二行使用另一个三角恒等式再施加 \(F\)。对于余结合律，将单位的自然性用于 \(\eta_{GY}:GY\rightarrow GFGY\)，得到
\[
GF(\eta_{GY})\circ\eta_{GY}
=\eta_{GFGY}\circ\eta_{GY}.
\]
施加 \(F\)，两边分别是 \(Q(\delta_Y)\circ\delta_Y\) 与 \(\delta_{QY}\circ\delta_Y\)，故余结合律也成立。这些核验与在相反范畴中应用\cref{b4:prop:adjunction-monad}得到的对偶证明一致。

\ProseParagraphBreak
单子记录可以复合的运算，而由伴随得到它时，另一范畴的全部信息未必都被记录。\Cref{b4:thm:A-em-adjunction}构造单子代数的自由与遗忘伴随；\cref{b4:thm:A-beck-monadicity}再判断原范畴何时能从这些代数恢复。它要求右伴随保守，即右伴随把一个态射送成同构时，该态射本身也必须是同构。另一个条件涉及分裂余等化子：它在余等化图上带有满足特定恒等式的截面与收缩映射，准确条件见\cref{b4:def:A-split-coequalizer}。判据要求每对经右伴随后具有这种分裂图的平行态射，在原范畴中都有余等化子，并被右伴随保持；不能把它理解为伴随自动保持所有余等化子。这是另外的构造与判据，并非伴随本身已经保证的结论。


\begin{chaptersummary}
可表函子把一个对象的泛性质组织成随目标或源自然变化的双射。表示对象必须连同这个双射或泛元素一起比较；保持这些数据的同构才是唯一确定的。Yoneda 引理把自然变换的全部分量缩减为在恒等态射上的取值，并说明这种恢复在对象和函子两个参数上仍然自然。

\ProseParagraphBreak
伴随把两个范畴之间的 Hom 集联系起来。单位与余单位来自恒等态射，转置公式说明它们已经决定全部双射，三角恒等式则恰好保证正反两种转置互逆。左伴随保持余极限、右伴随保持极限，都是把锥的相容条件和唯一分解经自然双射传送的结果；未获保证的那一类保持性质仍需另外检验。

\ProseParagraphBreak
单子保留自函子上的单位与归并运算，余单子反转这些箭头。自由模的嵌套形式和给出了可直接计算的例子。它们在这里解释伴随所附带的结构；怎样从单子恢复一个代数范畴，仍是另一个问题。
\end{chaptersummary}

\BookExercises

\begin{exercise}\label{b4:ex:03-torsion-representations}
设 \(m,n\ge1\)，对交换群 \(A\) 定义 \(T_n(A)=\{a\in A:na=0\}\)，并以群同态的限制定义态射作用。证明 \(T_n:\mathbf{Ab}\rightarrow\mathbf{Set}\) 可表。求从 \(T_m\) 到 \(T_n\) 的全部自然变换，并具体列出 \(m=6,n=4\) 时的答案。
\end{exercise}
\begin{solution}
同态 \(\mathbb Z/n\mathbb Z\rightarrow A\) 由 \(\bar1\) 的像唯一决定，该像恰为一个被 \(n\) 杀死的元素。因此 \(T_n\cong h^{\mathbb Z/n\mathbb Z}\)，且该对应与后复合相容。同样有 \(T_m\cong h^{\mathbb Z/m\mathbb Z}\)。协变 Yoneda 公式为 \(\operatorname{Nat}(h^A,h^B)\cong\operatorname{Hom}(B,A)\)：一个 \(s:B\rightarrow A\) 通过前复合，将每个 \(f:A\rightarrow X\) 变成 \(fs:B\rightarrow X\)。所以本题表示对象的态射从 \(\mathbb Z/n\mathbb Z\) 指向 \(\mathbb Z/m\mathbb Z\)，得到
\[
\operatorname{Nat}(T_m,T_n)
\cong\operatorname{Hom}_{\mathbb Z}(\mathbb Z/n\mathbb Z,\mathbb Z/m\mathbb Z).
\]
右边由满足 \(nc\equiv0\pmod m\) 的剩余类 \(\bar c\) 参数化，对应变换为 \(a\mapsto ca\)。这个公式不依赖 \(c\) 的代表，因为 \(ma=0\)；条件保证 \(nca=0\)。可选剩余类共有 \(\gcd(m,n)\) 个。若 \(m=6,n=4\)，条件为 \(6\mid4c\)，故 \(c\equiv0,3\pmod6\)，两变换分别为零与乘三。
\end{solution}

\begin{exercise}\label{b4:ex:03-units-nilpotents}
在允许零环的含幺交换环及保幺同态的范畴中，证明可逆元函子 \(A\mapsto A^\times\) 与函子 \(A\mapsto\{a\in A:a^n=0\}\)（固定 \(n\ge1\)）都可表，态射作用均为环同态的限制。分别写出表示对象及泛元素。再令
\[
\operatorname{Nil}(A)=\{a\in A:\;\text{存在整数}\;r\ge1\text{，使}\;a^r=0\},
\]
并以同样方式定义态射作用。证明全部幂零元函子 \(\operatorname{Nil}\) 不可表。
\end{exercise}
\begin{solution}
可逆元函子由 \(\mathbb Z[t,t^{-1}]\) 表示，泛元素为 \(t\)：给定可逆元 \(a\)，唯一同态为 \(t\mapsto a,t^{-1}\mapsto a^{-1}\)。关系 \(tt^{-1}=1\) 保证每个同态确实给出可逆元。第二个函子由 \(\mathbb Z[t]/(t^n)\) 表示，泛元素为 \(\bar t\)，因为指定 \(t\mapsto a\) 能通过商环当且仅当 \(a^n=0\)。两种对应均与环同态后复合相容。

环同态保持幂零性，所以 \(\operatorname{Nil}\) 确实是函子。由\cref{b4:prop:representables-preserve-limits}，若它可表，就必须保持所有小积。取一个域 \(k\)，令 \(A_r=k[t]/(t^r)\)（\(r\ge1\)），记 \(\bar t_r\) 为 \(t\) 在 \(A_r\) 中的像。每个 \(\bar t_r\) 都幂零，故 \((\bar t_r)_r\) 属于 \(\prod_{r\ge1}\operatorname{Nil}(A_r)\)。然而对任意固定整数 \(q\ge1\)，只要 \(r>q\)，就有 \(\bar t_r^{\,q}\ne0\)。因此这个元组在积环 \(\prod_{r\ge1}A_r\) 中不幂零，投影诱导的比较映射
\[
\operatorname{Nil}\!\left(\prod_{r\ge1}A_r\right)
\longrightarrow\prod_{r\ge1}\operatorname{Nil}(A_r)
\]
不为满射。这与保持积矛盾，所以 \(\operatorname{Nil}\) 不可表。
\end{solution}

\begin{exercise}\label{b4:ex:03-constant-representable}
设 \(\mathcal C\) 为局部小范畴。证明常值单元素协变函子可表，当且仅当 \(\mathcal C\) 有始对象；反变版本对应终对象。再证明常值空集合函子在任意 \(\mathcal C\) 上都不可表。
\end{exercise}
\begin{solution}
表示对象 \(A\) 满足 \(\operatorname{Hom}_{\mathcal C}(A,X)\) 对每个 \(X\) 恰有一个元素，这正是始对象条件。所有到单元素集合的映射自动相容，所以不存在额外自然性要求。反变版本同样把条件写成每个 \(\operatorname{Hom}_{\mathcal C}(X,A)\) 为单元素集，得到终对象。若空集合函子由 \(A\) 表示，则在 \(A\) 处得到 \(\operatorname{Hom}_{\mathcal C}(A,A)\cong\varnothing\)，与恒等态射存在矛盾；若 \(\mathcal C\) 没有对象，更没有可选的表示对象。
\end{solution}

\begin{exercise}\label{b4:ex:03-powerset-operations}
令 \(\mathcal P:\mathbf{Set}^{\mathrm{op}}\rightarrow\mathbf{Set}\) 为以逆像定义态射作用的幂集函子。求 \(\mathcal P\) 的全部自然自变换，不事先要求它们保持并、交或补。
\end{exercise}
\begin{solution}
由特征函数可知 \(\mathcal P\cong h_{\{0,1\}}\)。Yoneda 引理把自然自变换对应到四个映射 \(c:\{0,1\}\rightarrow\{0,1\}\)，变换作用为 \(\chi_A\mapsto c\circ\chi_A\)。恒等、交换 \(0,1\)、常值零、常值一，分别给出
\[
A\longmapsto A,\qquad A\longmapsto X\setminus A,
\qquad A\longmapsto\varnothing,\qquad A\longmapsto X
\quad(A\subseteq X).
\]
这些操作均与逆像交换。由 Yoneda 引理的双射，全部自然自变换恰好对应 \(\{0,1\}\) 的这四个集合自映射，因此没有其他可能。
\end{solution}

\begin{exercise}\label{b4:ex:03-natural-multioperations}
固定非负整数 \(r\)。对每个交换群 \(A\)，给定集合映射 \(t_A:A^r\rightarrow A\)，并要求其对群同态自然。证明存在唯一整数族 \((n_1,\ldots,n_r)\)，使
\(t_A(a_1,\ldots,a_r)=\sum_i n_i a_i\)。说明 \(r=0\) 的含义。
\end{exercise}
\begin{solution}
\(A\mapsto A^r\) 由带标准基 \(e_1,\ldots,e_r\) 的 \(\mathbb Z^r\) 表示。令 \(v=t_{\mathbb Z^r}(e_1,\ldots,e_r)=\sum_i n_ie_i\)。对唯一满足 \(f(e_i)=a_i\) 的群同态 \(f:\mathbb Z^r\rightarrow A\)，自然性给出 \(t_A(a_1,\ldots,a_r)=f(v)=\sum_i n_ia_i\)。取 \(A=\mathbb Z^r\) 及标准基恢复系数，故唯一；反过来这些整数线性组合均自然。\(r=0\) 时定义域是单元素集合，表示对象为零群，自然选择的元素只能是各群的零元。
\end{solution}

\begin{exercise}\label{b4:ex:03-center-natural-endomorphisms}
设 \(R\) 为含幺环，\(\mathcal M=R\text{-}\mathbf{Mod}\) 为幺性左 \(R\) 模范畴，\(U:\mathcal M\rightarrow\mathbf{Set}\) 为忘却函子。证明 \(\operatorname{Nat}(1_{\mathcal M},1_{\mathcal M})\) 与环中心 \(Z(R)\) 一一对应；再证明 \(\operatorname{Nat}(U,U)\) 作为集合与 \(R\) 一一对应。
\end{exercise}
\begin{solution}
先计算忘却函子的自然自变换。设 \(\alpha:U\Rightarrow U\)，令 \(a=\alpha_R(1)\)。对左线性映射 \(f_m:R\rightarrow M\)、\(r\mapsto rm\)，自然性给出 \(\alpha_M(m)=f_m(a)=am\)。反之，固定 \(a\in R\) 后，对任意模同态 \(f:M\rightarrow N\)，都有 \(f(am)=af(m)\)，故左乘 \(a\) 的集合映射组成 \(U\) 的自然自变换。取 \(M=R,m=1\) 恢复参数 \(a\)，所以对应唯一。

恒等函子 \(1_{\mathcal M}\) 的自然自变换还要求每个分量为 \(R\) 模同态。其底集合映射由上面的计算写成左乘 \(a\)；在 \(R\) 上，线性性要求 \(ar=ra\) 对所有 \(r\in R\) 成立，即 \(a\in Z(R)\)。反之，中心元素的左乘在每个左模上确实左线性，所以恰给出 \(1_{\mathcal M}\) 的全部自然自变换。自然性与分量的线性在这里是两个不同条件。
\end{solution}

\begin{exercise}\label{b4:ex:03-adjoint-unique}
设局部小范畴之间的函子 \(F:\mathcal C\rightarrow\mathcal D\) 同时以 \(G,G':\mathcal D\rightarrow\mathcal C\) 为右伴随，伴随双射已指定。证明存在唯一与这两组双射相容的自然同构 \(G\cong G'\)。
\end{exercise}
\begin{solution}
分别记 \(F\dashv G\) 与 \(F\dashv G'\) 的指定伴随双射为 \(\Phi,\Phi'\)。对每个 \(Y\in\mathcal D\)，它们复合给出自然于 \(X\) 的同构
\[
\operatorname{Hom}_{\mathcal C}(X,GY)
\xrightarrow{\Phi^{-1}}\operatorname{Hom}_{\mathcal D}(FX,Y)
\xrightarrow{\Phi'}\operatorname{Hom}_{\mathcal C}(X,G'Y).
\]
Yoneda 引理给出唯一同构 \(s_Y:GY\rightarrow G'Y\)，使上述映射为后复合 \(s_Y\)。若 \(b:Y\rightarrow Y'\)，两组伴随在 \(Y\) 上自然，故后复合 \(G'(b)s_Y\) 与后复合 \(s_{Y'}G(b)\) 对所有 \(X\) 相同。Yoneda 忠实性给出 \(G'(b)s_Y=s_{Y'}G(b)\)，所以 \(s\) 自然。逆双射同样产生自然逆，且相容性逐个 \(Y\) 确定 \(s_Y\)，故唯一。
\end{solution}

\begin{exercise}\label{b4:ex:03-composite-adjunction}
设 \(F:\mathcal C\rightarrow\mathcal D\)、\(G:\mathcal D\rightarrow\mathcal C\) 与 \(F':\mathcal D\rightarrow\mathcal E\)、\(G':\mathcal E\rightarrow\mathcal D\) 是局部小范畴之间的伴随 \(F\dashv G\)、\(F'\dashv G'\)。证明 \(F'F\dashv GG'\)，并用两组单位、余单位写出复合伴随的单位与余单位。
\end{exercise}
\begin{solution}
按顺序复合自然双射
\[
\operatorname{Hom}_{\mathcal E}(F'FX,Z)
\cong\operatorname{Hom}_{\mathcal D}(FX,G'Z)
\cong\operatorname{Hom}_{\mathcal C}(X,GG'Z).
\]
复合仍对 \(X,Z\) 自然，故给出伴随。将 \(1_{F'FX}\) 代入，先得到 \(\eta'_{FX}\)，再得到 \(G(\eta'_{FX})\eta_X\)，所以单位为该复合。将 \(1_{GG'Z}\) 经两个逆双射送回，得到余单位
\[
F'FGG'Z\xrightarrow{F'(\varepsilon_{G'Z})}F'G'Z
\xrightarrow{\varepsilon'_Z}Z.
\]
这些映射来自所构造的伴随双射，因而由\cref{b4:thm:adjunction-unit-counit}满足三角恒等式。
\end{solution}

\begin{exercise}\label{b4:ex:03-rounding-adjoints}
把 \(\mathbb Z\) 与 \(\mathbb R\) 按通常次序看成范畴，令 \(i:\mathbb Z\rightarrow\mathbb R\) 为包含。证明 \(\lceil-\rceil\dashv i\dashv\lfloor-\rfloor\)，并解释单位、余单位是什么。
\end{exercise}
\begin{solution}
偏序范畴的每个 Hom 集为空或单元素，所以伴随双射等价于相应不等式同时成立。对 \(x\in\mathbb R,n\in\mathbb Z\)，
\[
\lceil x\rceil\le n\iff x\le n,
\qquad
n\le x\iff n\le\lfloor x\rfloor.
\]
向上取整、向下取整均单调，故为函子；偏序范畴的任意两个具有相同源和目标的态射自动相等，因此双射自然。\(\lceil-\rceil\dashv i\) 的单位为 \(x\le\lceil x\rceil\)，余单位为整数处的等式 \(\lceil n\rceil=n\)；\(i\dashv\lfloor-\rfloor\) 的单位为 \(n=\lfloor n\rfloor\)，余单位为 \(\lfloor x\rfloor\le x\)。
\end{solution}

\begin{exercise}\label{b4:ex:03-set-exponential-adjunction}
固定集合 \(A\)。证明 \(-\times A:\mathbf{Set}\rightarrow\mathbf{Set}\) 左伴随于 \((-)^A\)，写出单位、余单位，并直接验证三角恒等式。包括 \(A=\varnothing\) 的情形。
\end{exercise}
\begin{solution}
双射把 \(f:X\times A\rightarrow Y\) 送到 \(x\mapsto(a\mapsto f(x,a))\)，逆把 \(g:X\rightarrow Y^A\) 送到 \((x,a)\mapsto g(x)(a)\)。前复合 \(u:X'\rightarrow X\) 或后复合 \(v:Y\rightarrow Y'\) 后，两条路径在 \((x',a)\) 上均取值 \(v(f(u(x'),a))\)，故自然。单位为 \(\eta_X(x)(a)=(x,a)\)，余单位为求值 \(\varepsilon_Y(h,a)=h(a)\)。第一条三角恒等式将 \((x,a)\) 送到 \(\eta_X(x)(a)=(x,a)\)；第二条将 \(h:A\rightarrow Y\) 送到函数 \(a\mapsto h(a)\)。若 \(A\) 空，\(\eta_X(x)\) 就是从空集到 \(X\times\varnothing=\varnothing\) 的唯一函数；\(Y^\varnothing\) 也只有从空集到 \(Y\) 的唯一函数。此时左函子常值为空集，右函子常值为单元素集，余单位为从空集到 \(Y\) 的唯一函数，两边的 Hom 集均为单元素集，三角恒等式也比较相同源、目标之间的唯一函数。
\end{solution}

\begin{exercise}\label{b4:ex:03-s3-abelianization}
求 \(S_3\) 的交换化，并据此求对任意交换群 \(A\) 的全部群同态 \(S_3\rightarrow A\)。把结果与伴随的单位联系起来。
\end{exercise}
\begin{solution}
符号同态的目标交换，故 \([S_3,S_3]\subseteq A_3\)。按交换子约定 \([x,y]=xyx^{-1}y^{-1}\)，\([(12),(23)]=(132)\)，它生成 \(A_3\)，所以 \([S_3,S_3]=A_3\)。因此 \(S_3^{\mathrm{ab}}\cong C_2\)，单位就是符号商映射。到交换群 \(A\) 的每个同态唯一经过它，继而由 \(C_2\) 的生成元的像决定；该像是满足 \(2a=0\) 的任意元素 \(a\in A\)。具体地，偶置换映到零，奇置换映到 \(a\)。非平凡的 \(A_3\) 被单位消去，显示此伴随的单位并非同构。
\end{solution}

\begin{exercise}\label{b4:ex:03-fully-faithful-adjoint}
设 \(F:\mathcal C\rightarrow\mathcal D\) 左伴随于 \(G:\mathcal D\rightarrow\mathcal C\)，两范畴局部小。证明 \(G\) 全忠实，当且仅当每个余单位 \(\varepsilon_Y:FGY\rightarrow Y\) 是同构。
\end{exercise}
\begin{solution}
若余单位均为同构，则对 \(Y,Y'\in\mathcal D\)，复合双射
\[
\operatorname{Hom}_{\mathcal D}(Y,Y')
\xrightarrow{-\circ\varepsilon_Y}
\operatorname{Hom}_{\mathcal D}(FGY,Y')
\xrightarrow{\Phi}
\operatorname{Hom}_{\mathcal C}(GY,GY')
\]
把 \(h\) 送到 \(G(h)G(\varepsilon_Y)\eta_{GY}=G(h)\)。所以 \(G\) 在 Hom 集上给出双射。

反过来，设 \(G\) 全忠实。单位 \(\eta_{GY}:GY\rightarrow GFGY\) 属于 Hom 集映射
\[
G:\operatorname{Hom}_{\mathcal D}(Y,FGY)
\longrightarrow\operatorname{Hom}_{\mathcal C}(GY,GFGY)
\]
的目标。由 \(G\) 的全性，该集合映射满射，故存在 \(s_Y:Y\rightarrow FGY\) 使 \(G(s_Y)=\eta_{GY}\)。三角恒等式给出 \(G(\varepsilon_Ys_Y)=1_{GY}\)，忠实性推出 \(\varepsilon_Ys_Y=1_Y\)。另一方面，
\[
\Phi(s_Y\varepsilon_Y)
=G(s_Y)G(\varepsilon_Y)\eta_{GY}
=\eta_{GY}=\Phi(1_{FGY}).
\]
\(\Phi\) 的单射性给出 \(s_Y\varepsilon_Y=1_{FGY}\)，故 \(\varepsilon_Y\) 可逆。
\end{solution}

\begin{exercise}\label{b4:ex:03-torsion-quotient-adjunction}
固定 \(n>1\)，在交换群范畴中取 \(P=\mathbb Z/n\mathbb Z\)。把 Hom–张量伴随写成 \(M/nM\) 与 \(N[n]=\{x\in N:nx=0\}\) 的形式，并具体写出单位、余单位。对 \(M=N=\mathbb Z\) 分别计算它们。
\end{exercise}
\begin{solution}
自然同构为
\[
P\otimes_{\mathbb Z}M\cong M/nM,
\quad\bar1\otimes m\longmapsto\bar m,
\qquad
\operatorname{Hom}_{\mathbb Z}(P,N)\cong N[n],
\quad f\longmapsto f(\bar1).
\]
第一个逆为 \(\bar m\mapsto\bar1\otimes m\)，关系 \(\bar1\otimes nm=n\bar1\otimes m=0\) 保证良定义；第二个逆由生成元像唯一延拓。因此伴随为 \(M\mapsto M/nM\) 左伴随于 \(N\mapsto N[n]\)。因为 \(n\) 消去整个商群 \(M/nM\)，有 \((M/nM)[n]=M/nM\)，单位便是商映射 \(M\rightarrow M/nM\)。又因 \(n\) 消去 \(N[n]\)，有 \(nN[n]=0\)，故 \(N[n]/nN[n]=N[n]\)，余单位便是它到 \(N\) 的包含。对 \(\mathbb Z\)，前者为 \(\mathbb Z\rightarrow\mathbb Z/n\mathbb Z\)，后者为 \(0\rightarrow\mathbb Z\)，两者均非同构。
\end{solution}

\begin{exercise}\label{b4:ex:03-tensor-equalizer-failure}
在交换群范畴中取平行态射 \(2,0:\mathbb Z\rightrightarrows\mathbb Z\)。证明左伴随 \((\mathbb Z/2\mathbb Z)\otimes_{\mathbb Z}-\) 不保持这对态射的等化子。
\end{exercise}
\begin{solution}
原等化子为 \(\{m\in\mathbb Z:2m=0\}=0\)。张量后两个态射都成为 \(\mathbb Z/2\mathbb Z\) 上的零映射，其等化子为整个 \(\mathbb Z/2\mathbb Z\)。原等化子的张量仍为零，比较映射是 \(0\rightarrow\mathbb Z/2\mathbb Z\)，不是同构。这表明一般左伴随甚至不必保持有限极限，尽管它保持所有存在的小余极限。
\end{solution}

\begin{exercise}\label{b4:ex:03-free-product-comparison}
令 \(F:\mathbf{Set}\rightarrow\mathbf{Ab}\) 为自由交换群函子。计算 \(X=Y=\{*\}\) 时的积比较映射 \(F(X\times Y)\rightarrow F(X)\times F(Y)\)，并判断 \(F\) 是否保持终对象及二元积。
\end{exercise}
\begin{solution}
源为 \(\mathbb Z\)，目标为 \(\mathbb Z^2\)。两个投影都把唯一生成元送到相应的生成元，所以比较映射为 \(m\mapsto(m,m)\)，并不满射。因此 \(F\) 不保持二元积。集合范畴的终对象是单元素集合，它的像为 \(\mathbb Z\)，而交换群范畴的终对象为零群，所以也不保持终对象。这里的自由基来自指定集合，问题不在于选择了哪一组基。
\end{solution}

\begin{exercise}\label{b4:ex:03-constant-adjoints}
固定集合 \(A\)，令 \(K_A:\mathbf{Set}\rightarrow\mathbf{Set}\) 为常值 \(A\) 的函子，所有态射映成 \(1_A\)。分别判断它什么时候有右伴随、什么时候有左伴随，并写出相应伴随。
\end{exercise}
\begin{solution}
若 \(K_A\) 有右伴随，它是左伴随，必须保持始对象 \(\varnothing\)，故 \(A=\varnothing\)。反过来，\(\operatorname{Hom}(\varnothing,Y)\) 与 \(\operatorname{Hom}(X,\{*\})\) 恒为单元素集，给出自然双射，故 \(K_\varnothing\dashv K_{\{*\}}\)。若 \(K_A\) 有左伴随，它是右伴随，必须保持终对象，故 \(A\) 为单元素集合；刚才的伴随给出充分性。于是这两类情形完全确定。
\end{solution}

\begin{exercise}\label{b4:ex:03-torsion-inverse-system}
固定素数 \(p\)，以约化映射组成逆系统 \(\mathbb Z/p^r\mathbb Z\)，\(r\ge1\)，令 \(L\) 为其逆极限。证明每个 \((\mathbb Z/p^r\mathbb Z)[p]\) 都非零，而 \(L[p]=0\)。解释这为何不违背右伴随保持逆极限。
\end{exercise}
\begin{solution}
第 \(r\) 层被 \(p\) 杀死的元素为 \(\bar c p^{r-1}\)，\(c\in\mathbb Z/p\mathbb Z\)，所以组成阶为 \(p\) 的非零群。但从第 \(r+1\) 层到第 \(r\) 层的限制映射把 \(\bar c p^r\) 送到零，因此这些 \(p\) 挠子群构成的逆系统，其转移映射全为零。任何相容族 \((x_r)_r\) 都必须满足 \(x_r=0(x_{r+1})=0\)，只能是零族。

直接地，若 \((a_r)_r\in L\) 被 \(p\) 杀死，则 \(a_{r+1}\) 在第 \(r+1\) 层为 \(p^r\) 的倍数，它约化到第 \(r\) 层为零，故每个 \(a_r=0\)。函子 \(M\mapsto M[p]\cong\operatorname{Hom}_{\mathbb Z}(\mathbb Z/p\mathbb Z,M)\) 是右伴随，所得同构恰为
\[
L[p]\cong\varprojlim_r(\mathbb Z/p^r\mathbb Z)[p]=0.
\]
右伴随所保持的是相容族的极限，上式两边均为零；逐层非零性并不保证存在非零相容族。
\end{solution}

\begin{exercise}\label{b4:ex:03-reader-monad}
固定集合 \(A\)，令 \(TX=X^A\)，\(T(f)\) 为后复合 \(f\)。定义 \(\eta_X(x)\) 为常值 \(x\) 的函数，并对 \(u:A\rightarrow X^A\) 定义 \(\mu_X(u)(a)=u(a)(a)\)。证明 \((T,\eta,\mu)\) 是集合范畴上的单子。
\end{exercise}
\begin{solution}
\(T\) 的恒等与复合律来自函数复合。对 \(f:X\rightarrow Y\)，两种把 \(u\in T^2X\) 送到 \(TY\) 的方式在 \(a\) 处都取值 \(f(u(a)(a))\)，所以 \(\mu\) 自然；\(\eta\) 的自然性来自常值函数的后复合。对 \(v\in TX\)，\(T\eta_X(v)\) 在 \(a\) 处是常值 \(v(a)\) 的函数，故 \(\mu_XT\eta_X(v)=v\)；\(\eta_{TX}(v)\) 为常值 \(v\) 的函数，故 \(\mu_X\eta_{TX}(v)=v\)。最后，对 \(w\in T^3X\)，两种归并在 \(a\) 处都得到 \(w(a)(a)(a)\)，从而结合律成立。若 \(A\) 空，各函数集合为单元素集，以上等式仍成立。
\end{solution}

\begin{exercise}\label{b4:ex:03-product-comonad}
固定集合 \(A\)，令 \(QX=A\times X\)，\(Qf=1_A\times f\)。定义 \(\varepsilon_X(a,x)=x\)、\(\delta_X(a,x)=(a,(a,x))\)。证明这给出余单子；若 \(A=\varnothing\)，结论是否仍成立？
\end{exercise}
\begin{solution}
对任意函数 \(f\)，先重复 \(A\) 坐标再应用 \(f\)，与先应用 \(f\) 再重复坐标结果相同，所以 \(\delta\) 自然；投影的自然性由 \(f(\varepsilon_X(a,x))=f(x)=\varepsilon_Y(a,f(x))\) 给出。两种余结合复合都把 \((a,x)\) 送到 \((a,(a,(a,x)))\)。复合 \(Q\varepsilon_X\circ\delta_X\) 与 \(\varepsilon_{QX}\circ\delta_X\) 都把它送回 \((a,x)\)。所以三条公理成立。\(A\) 空时 \(Q\) 常值为空集合，所有公理比较的映射都有空定义域，仍成立。
\end{solution}

\begin{exercise}\label{b4:ex:03-list-monad}
对集合 \(X\)，令 \(TX=\coprod_{n\ge0}X^n\) 为全部有限有序表的集合，包括空表；对函数逐项应用。以单元素表定义 \(\eta_X:X\rightarrow TX\)，以依次连接各子表定义 \(\mu_X:T^2X\rightarrow TX\)。证明这是一个单子，并说明空表在这个构造中的作用。
\end{exercise}
\begin{solution}
逐项应用函数保持表长和次序，因此给出函子。逐项应用函数与插入单元素表、连接各子表相容，所以 \(\eta,\mu\) 自然。对表 \((x_1,\ldots,x_n)\)，先把每项换成单元素表再连接，或先把整个表视为唯一子表再连接，都恢复原表，这给出两条单位律。对三层表，两种归并都按先外层位置、再中层位置、最后内层位置的次序输出同一列元素，故结合律成立；空子表贡献零个元素，不影响论证。

空表使 \(T\varnothing\) 仍有一个元素，并允许连接中出现零个子表；在将此构造解释为自由幺半群时，它正是单位元。若全部改为非空有限表，也会得到另一个满足单子公理的构造，但它表达的是不要求单位元的半群，已不是本题所定义的全部有限表。
\end{solution}
